% programmes-cpge-mpsi.tex — programme de physique-chimie de la classe de MPSI
% Texte officiel de 2020 (voie Mathématiques, physique et sciences de l’ingénieur), publié par le ministère de
% l’Enseignement supérieur. Fichier de données du paquet programmes-cpge.
% Format : voir la documentation, « Ajouter un programme ».

\cpgefiliere{MPSI}{
  discipline = {Physique-chimie},
  classe = {MPSI},
  voie = {Mathématiques, physique et sciences de l’ingénieur},
  annee = {2020},
}

\cpgerubrique{MPSI/P}{
  parent = {MPSI},
  numero = {P},
  niveau = {1},
  titre = {Préambule},
  chapeau = {},
}

\cpgerubrique{MPSI/P/1}{
  parent = {MPSI/P},
  numero = {P.1},
  niveau = {2},
  titre = {Objectifs de formation},
  chapeau = {
    \cpgeitem{Le programme de physique-chimie de la classe de MPSI est conçu comme un socle cohérent et ambitieux de connaissances et de capacités scientifiques préparant les étudiants à la deuxième année de classe préparatoire et, au-delà, à un cursus d’ingénieur, de chercheur ou d’enseignant. Il s’agit de renforcer chez l’étudiant les compétences déjà travaillées au lycée inhérentes à la pratique de la démarche scientifique : observer et s’approprier, analyser et modéliser, réaliser et valider, et enfin communiquer et valoriser ses résultats.}
    \cpgeitem{L’acquisition de ce socle par les étudiants constitue un objectif prioritaire pour l’enseignant. Parce que la physique et la chimie sont avant tout des sciences expérimentales qui développent la curiosité, la créativité et l’analyse critique, l’expérience est au cœur de son enseignement, que ce soit en cours ou lors des séances de travaux pratiques. Les activités expérimentales habituent les étudiants à se confronter au réel, comme ils auront à le faire dans l’exercice de leur métier.}
    \cpgeitem{De même, l’introduction de capacités numériques dans le programme prend en compte la place nouvelle des sciences numériques dans la formation des scientifiques notamment dans le domaine de la simulation. Elles offrent aux étudiants la possibilité d’effectuer une modélisation avancée du monde réel, par exemple par la prise en compte d’effets non linéaires.}
    \cpgeitem{La démarche de modélisation occupe également une place centrale dans le programme pour former les étudiants à établir, de manière autonome, un lien fait d’allers-retours entre le « monde » des objets, des expériences, des faits, et celui des modèles et des théories. L’enseignant doit rechercher un point d’équilibre entre des approches complémentaires : conceptuelle et expérimentale, abstraite et concrète, théorique et appliquée, inductive et déductive, qualitative et quantitative. La construction d’un modèle passe aussi par l’utilisation maîtrisée des mathématiques dont un des fondateurs de la physique expérimentale, Galilée, énonçait déjà qu’elles sont le langage dans lequel est écrit le monde.}
    \cpgeitem{Enfin, l’autonomie et la prise d’initiative sont spécifiquement développées à travers la pratique d’activités du type « résolution de problèmes » qui visent à exercer les étudiants à mobiliser de façon complémentaire connaissances et capacités pour répondre à un questionnement ou atteindre un but sans qu’aucune démarche de résolution ne soit fournie.}},
}

\cpgerubrique{MPSI/P/2}{
  parent = {MPSI/P},
  numero = {P.2},
  niveau = {2},
  titre = {Organisation du programme},
  chapeau = {
    \cpgeitem{Le programme est organisé en deux parties.}
    \cpgeitem{Dans la première partie, intitulée « Formation expérimentale », sont décrits les objectifs de formation sur le thème « Mesures et incertitudes » ainsi que les méthodes et les capacités expérimentales que les étudiants doivent maîtriser à la fin de l’année scolaire. Leur mise en œuvre doit notamment s’appuyer sur des problématiques concrètes identifiées en gras dans la seconde partie du programme intitulée « Contenus thématiques ». Elles doivent être programmées par l’enseignant de façon à assurer un apprentissage progressif de l’ensemble des capacités attendues.}
    \cpgeitem{La seconde partie, intitulée « Contenus thématiques » est structurée autour de quatre thèmes : « ondes et signaux », « mouvements et interactions », « l’énergie : conversions et transferts » et « constitution et transformations de la matière ». La présentation en deux colonnes (« notions et contenus » et « capacités exigibles ») met en valeur les éléments clefs constituant le socle de connaissances et de capacités dont l’assimilation par tous les étudiants est requise. La progression dans les contenus disciplinaires est organisée en deux semestres. Pour faciliter la progressivité des acquisitions, au premier semestre les grandeurs physiques introduites sont essentiellement des grandeurs scalaires dépendant du temps et éventuellement d’une variable d’espace. Certains items de cette seconde partie, identifiés en caractères gras, se prêtent particulièrement à une approche expérimentale. Ils doivent être abordés en priorité lors de séances de travaux pratiques où l’autonomie et l’initiative individuelle de l’étudiant doivent être privilégiées. La présence de capacités numériques explicitées atteste par ailleurs de la volonté de renforcer ce volet de la formation des étudiants.}
    \cpgeitem{Trois annexes sont consacrées d’une part au matériel nécessaire à la mise en œuvre des programmes, d’autre part aux outils mathématiques et aux outils numériques que les étudiants doivent savoir mobiliser de façon autonome dans le cadre des enseignements de physique-chimie à la fin de l’année de la classe de MPSI.}
    \cpgeitem{Ce programme précise les objectifs de formation à atteindre pour tous les étudiants. Il n’impose en aucun cas une progression pour chacun des deux semestres ; celle-ci relève de la liberté pédagogique de l’enseignant.}},
}

\cpgerubrique{MPSI/P/3}{
  parent = {MPSI/P},
  numero = {P.3},
  niveau = {2},
  titre = {Les compétences travaillées dans le cadre de la démarche scientifique},
  chapeau = {
    \cpgeitem{L’ensemble des activités proposées en classe préparatoire aux grandes écoles – activités expérimentales, résolutions de problèmes, TIPE, etc. – permet de travailler les compétences de la démarche scientifique qui figurent dans le tableau ci-dessous. Des capacités associées sont explicitées afin de préciser les contours de chaque compétence, elles ne constituent donc pas une liste exhaustive et peuvent parfois relever de plusieurs domaines de compétences. L’ordre de présentation de ces compétences ne préjuge pas d’un ordre de mobilisation de ces dernières lors d’une activité.}
    \cpgeitem{Les compétences doivent être acquises à l’issue de la formation en CPGE. Elles nécessitent d’être régulièrement mobilisées par les étudiants et sont évaluées en s’appuyant, par exemple, sur l’utilisation de grilles d’évaluation.}
    \cpgeitem{Le niveau de maîtrise de ces compétences dépend de \textbf{l’autonomie et de l’initiative} requises dans les activités proposées aux étudiants sur les notions et capacités exigibles du programme.}
    \cpgeitem{La mise en œuvre des programmes doit aussi être l’occasion d’aborder avec les étudiants des questions liées à l’histoire de l’évolution des idées, des modèles et des théories en physique-chimie, à des questions liées à la recherche scientifique actuelle et à des enjeux citoyens comme la responsabilité individuelle et collective, la \textbf{sécurité} pour soi et pour autrui, l’\textbf{environnement} et le \textbf{développement durable} ou encore le \textbf{réchauffement climatique}.}},
}

\cpgerubrique{MPSI/P/3/APP}{
  parent = {MPSI/P/3},
  numero = {APP},
  niveau = {3},
  titre = {S’approprier},
  chapeau = {},
}

\cpgeligne{MPSI/P/3/APP/1}{
  parent = {MPSI/P/3/APP},
  numero = {APP.1},
  contenu = {
    \cpgeitem{S’approprier}},
  capacites = {
    \cpgeitem{Rechercher, extraire et organiser l’information en lien avec la situation étudiée.}},
  experimentales = {},
  numeriques = {},
}

\cpgeligne{MPSI/P/3/APP/2}{
  parent = {MPSI/P/3/APP},
  numero = {APP.2},
  contenu = {
    \cpgeitem{S’approprier}},
  capacites = {
    \cpgeitem{Identifier la complémentarité d’informations présentées sous des formes différentes (texte, graphe, tableau, etc.).}},
  experimentales = {},
  numeriques = {},
}

\cpgeligne{MPSI/P/3/APP/3}{
  parent = {MPSI/P/3/APP},
  numero = {APP.3},
  contenu = {
    \cpgeitem{S’approprier}},
  capacites = {
    \cpgeitem{Énoncer ou dégager une problématique scientifique.}},
  experimentales = {},
  numeriques = {},
}

\cpgeligne{MPSI/P/3/APP/4}{
  parent = {MPSI/P/3/APP},
  numero = {APP.4},
  contenu = {
    \cpgeitem{S’approprier}},
  capacites = {
    \cpgeitem{Représenter la situation par un schéma modèle.}},
  experimentales = {},
  numeriques = {},
}

\cpgeligne{MPSI/P/3/APP/5}{
  parent = {MPSI/P/3/APP},
  numero = {APP.5},
  contenu = {
    \cpgeitem{S’approprier}},
  capacites = {
    \cpgeitem{Identifier les grandeurs pertinentes, leur attribuer un symbole.}},
  experimentales = {},
  numeriques = {},
}

\cpgeligne{MPSI/P/3/APP/6}{
  parent = {MPSI/P/3/APP},
  numero = {APP.6},
  contenu = {
    \cpgeitem{S’approprier}},
  capacites = {
    \cpgeitem{Relier le problème à une situation modèle connue.}},
  experimentales = {},
  numeriques = {},
}

\cpgeligne{MPSI/P/3/APP/7}{
  parent = {MPSI/P/3/APP},
  numero = {APP.7},
  contenu = {
    \cpgeitem{S’approprier}},
  capacites = {
    \cpgeitem{Acquérir de nouvelles connaissances en autonomie.}},
  experimentales = {},
  numeriques = {},
}

\cpgerubrique{MPSI/P/3/ANA}{
  parent = {MPSI/P/3},
  numero = {ANA},
  niveau = {3},
  titre = {Analyser/Raisonner},
  chapeau = {},
}

\cpgeligne{MPSI/P/3/ANA/1}{
  parent = {MPSI/P/3/ANA},
  numero = {ANA.1},
  contenu = {
    \cpgeitem{Analyser / Raisonner}},
  capacites = {
    \cpgeitem{Formuler des hypothèses.}},
  experimentales = {},
  numeriques = {},
}

\cpgeligne{MPSI/P/3/ANA/2}{
  parent = {MPSI/P/3/ANA},
  numero = {ANA.2},
  contenu = {
    \cpgeitem{Analyser / Raisonner}},
  capacites = {
    \cpgeitem{Décomposer un problème en plusieurs problèmes plus simples.}},
  experimentales = {},
  numeriques = {},
}

\cpgeligne{MPSI/P/3/ANA/3}{
  parent = {MPSI/P/3/ANA},
  numero = {ANA.3},
  contenu = {
    \cpgeitem{Analyser / Raisonner}},
  capacites = {
    \cpgeitem{Proposer une stratégie pour répondre à une problématique.}},
  experimentales = {},
  numeriques = {},
}

\cpgeligne{MPSI/P/3/ANA/4}{
  parent = {MPSI/P/3/ANA},
  numero = {ANA.4},
  contenu = {
    \cpgeitem{Analyser / Raisonner}},
  capacites = {
    \cpgeitem{Choisir, concevoir, justifier un protocole, un dispositif expérimental, un modèle ou des lois physiques.}},
  experimentales = {},
  numeriques = {},
}

\cpgeligne{MPSI/P/3/ANA/5}{
  parent = {MPSI/P/3/ANA},
  numero = {ANA.5},
  contenu = {
    \cpgeitem{Analyser / Raisonner}},
  capacites = {
    \cpgeitem{Évaluer des ordres de grandeur.}},
  experimentales = {},
  numeriques = {},
}

\cpgeligne{MPSI/P/3/ANA/6}{
  parent = {MPSI/P/3/ANA},
  numero = {ANA.6},
  contenu = {
    \cpgeitem{Analyser / Raisonner}},
  capacites = {
    \cpgeitem{Identifier les idées essentielles d’un document et leurs articulations.}},
  experimentales = {},
  numeriques = {},
}

\cpgeligne{MPSI/P/3/ANA/7}{
  parent = {MPSI/P/3/ANA},
  numero = {ANA.7},
  contenu = {
    \cpgeitem{Analyser / Raisonner}},
  capacites = {
    \cpgeitem{Relier qualitativement ou quantitativement différents éléments d’un ou de documents.}},
  experimentales = {},
  numeriques = {},
}

\cpgerubrique{MPSI/P/3/REA}{
  parent = {MPSI/P/3},
  numero = {REA},
  niveau = {3},
  titre = {Réaliser},
  chapeau = {},
}

\cpgeligne{MPSI/P/3/REA/1}{
  parent = {MPSI/P/3/REA},
  numero = {REA.1},
  contenu = {
    \cpgeitem{Réaliser}},
  capacites = {
    \cpgeitem{Mettre en œuvre les étapes d’une démarche, un protocole, un modèle.}},
  experimentales = {},
  numeriques = {},
}

\cpgeligne{MPSI/P/3/REA/2}{
  parent = {MPSI/P/3/REA},
  numero = {REA.2},
  contenu = {
    \cpgeitem{Réaliser}},
  capacites = {
    \cpgeitem{Extraire une information d’un texte, d’un graphe, d’un tableau, d’un schéma, d’une photo.}},
  experimentales = {},
  numeriques = {},
}

\cpgeligne{MPSI/P/3/REA/3}{
  parent = {MPSI/P/3/REA},
  numero = {REA.3},
  contenu = {
    \cpgeitem{Réaliser}},
  capacites = {
    \cpgeitem{Schématiser un dispositif, une expérience, une méthode de mesure.}},
  experimentales = {},
  numeriques = {},
}

\cpgeligne{MPSI/P/3/REA/4}{
  parent = {MPSI/P/3/REA},
  numero = {REA.4},
  contenu = {
    \cpgeitem{Réaliser}},
  capacites = {
    \cpgeitem{Utiliser le matériel et les produits de manière adaptée en respectant des règles de sécurité.}},
  experimentales = {},
  numeriques = {},
}

\cpgeligne{MPSI/P/3/REA/5}{
  parent = {MPSI/P/3/REA},
  numero = {REA.5},
  contenu = {
    \cpgeitem{Réaliser}},
  capacites = {
    \cpgeitem{Effectuer des représentations graphiques à partir de données.}},
  experimentales = {},
  numeriques = {},
}

\cpgeligne{MPSI/P/3/REA/6}{
  parent = {MPSI/P/3/REA},
  numero = {REA.6},
  contenu = {
    \cpgeitem{Réaliser}},
  capacites = {
    \cpgeitem{Mener des calculs analytiques ou à l’aide d’un langage de programmation, effectuer des applications numériques.}},
  experimentales = {},
  numeriques = {},
}

\cpgeligne{MPSI/P/3/REA/7}{
  parent = {MPSI/P/3/REA},
  numero = {REA.7},
  contenu = {
    \cpgeitem{Réaliser}},
  capacites = {
    \cpgeitem{Conduire une analyse dimensionnelle.}},
  experimentales = {},
  numeriques = {},
}

\cpgerubrique{MPSI/P/3/VAL}{
  parent = {MPSI/P/3},
  numero = {VAL},
  niveau = {3},
  titre = {Valider},
  chapeau = {},
}

\cpgeligne{MPSI/P/3/VAL/1}{
  parent = {MPSI/P/3/VAL},
  numero = {VAL.1},
  contenu = {
    \cpgeitem{Valider}},
  capacites = {
    \cpgeitem{Exploiter des observations, des mesures en estimant les incertitudes.}},
  experimentales = {},
  numeriques = {},
}

\cpgeligne{MPSI/P/3/VAL/2}{
  parent = {MPSI/P/3/VAL},
  numero = {VAL.2},
  contenu = {
    \cpgeitem{Valider}},
  capacites = {
    \cpgeitem{Confronter les résultats d’un modèle à des résultats expérimentaux, à des données figurant dans un document, à ses connaissances.}},
  experimentales = {},
  numeriques = {},
}

\cpgeligne{MPSI/P/3/VAL/3}{
  parent = {MPSI/P/3/VAL},
  numero = {VAL.3},
  contenu = {
    \cpgeitem{Valider}},
  capacites = {
    \cpgeitem{Confirmer ou infirmer une hypothèse, une information.}},
  experimentales = {},
  numeriques = {},
}

\cpgeligne{MPSI/P/3/VAL/4}{
  parent = {MPSI/P/3/VAL},
  numero = {VAL.4},
  contenu = {
    \cpgeitem{Valider}},
  capacites = {
    \cpgeitem{Analyser les résultats de manière critique.}},
  experimentales = {},
  numeriques = {},
}

\cpgeligne{MPSI/P/3/VAL/5}{
  parent = {MPSI/P/3/VAL},
  numero = {VAL.5},
  contenu = {
    \cpgeitem{Valider}},
  capacites = {
    \cpgeitem{Repérer les points faibles d’une argumentation (contradiction, partialité, incomplétude, etc.).}},
  experimentales = {},
  numeriques = {},
}

\cpgeligne{MPSI/P/3/VAL/6}{
  parent = {MPSI/P/3/VAL},
  numero = {VAL.6},
  contenu = {
    \cpgeitem{Valider}},
  capacites = {
    \cpgeitem{Proposer des améliorations de la démarche ou du modèle.}},
  experimentales = {},
  numeriques = {},
}

\cpgerubrique{MPSI/P/3/COM}{
  parent = {MPSI/P/3},
  numero = {COM},
  niveau = {3},
  titre = {Communiquer},
  chapeau = {},
}

\cpgeligne{MPSI/P/3/COM/1}{
  parent = {MPSI/P/3/COM},
  numero = {COM.1},
  contenu = {
    \cpgeitem{Communiquer}},
  capacites = {
    \cpgeitem{À l’écrit comme à l’oral : \begin{itemize} \item présenter les étapes de sa démarche de manière synthétique, organisée et cohérente. \item rédiger une synthèse, une analyse, une argumentation. \item utiliser un vocabulaire scientifique précis et choisir des modes de représentation adaptés (schémas, graphes, cartes mentales, etc.). \end{itemize}}},
  experimentales = {},
  numeriques = {},
}

\cpgeligne{MPSI/P/3/COM/2}{
  parent = {MPSI/P/3/COM},
  numero = {COM.2},
  contenu = {
    \cpgeitem{Communiquer}},
  capacites = {
    \cpgeitem{Écouter, confronter son point de vue.}},
  experimentales = {},
  numeriques = {},
}

\cpgerubrique{MPSI/P/4}{
  parent = {MPSI/P},
  numero = {P.4},
  niveau = {2},
  titre = {Repères pour l’enseignement},
  chapeau = {
    \cpgeitem{Dans le cadre de la liberté pédagogique, l’enseignant organise son enseignement en respectant trois grands principes directeurs : \begin{itemize} \item privilégier la mise en activité des étudiants en évitant tout dogmatisme : l’acquisition des connaissances, des capacités et des compétences est d’autant plus efficace que les étudiants sont acteurs de leur formation. Les supports pédagogiques utilisés doivent notamment favoriser la réflexion, le raisonnement, la participation et l’autonomie des étudiants. L’investigation expérimentale et la résolution de problèmes favorisent cette mise en activité ; \item recourir à la mise en contexte des contenus scientifiques : le questionnement scientifique peut être introduit à partir de phénomènes naturels, de procédés industriels ou d’objets technologiques. Le recours à des approches documentaires est un moyen pertinent pour diversifier les supports d’accès à l’information scientifique et technologique et ainsi former l’étudiant à mieux en appréhender la complexité et à apprendre par lui-même. Lorsque le thème traité s’y prête, l’enseignant peut le mettre en perspective avec l’histoire des sciences et des techniques, avec des questions d’actualité ou des débats d’idées ; \item contribuer à la nécessaire mise en cohérence des enseignements scientifiques ; la progression en physique-chimie doit être articulée avec celles mises en œuvre dans les autres disciplines scientifiques : mathématiques, informatique, sciences industrielles de l’ingénieur. \end{itemize}}
    \cpgeitem{Concernant l’évaluation, qui vise à mesurer le degré de maîtrise du socle ainsi défini et le niveau d’autonomie et d’initiative des étudiants, l’enseignant veille soigneusement à identifier les compétences et les capacités mobilisées dans les activités proposées afin d’en élargir le plus possible le spectre.}
    \cpgeitem{Enfin, le professeur veille aussi à développer chez les étudiants des compétences transversales et préprofessionnelles relatives aux capacités suivantes : \begin{itemize} \item identifier les différents champs professionnels et les parcours pour y accéder ; \item valoriser ses compétences scientifiques et techniques en lien avec son projet de poursuite d'études ou professionnel. \end{itemize}}},
}

\cpgerubrique{MPSI/E}{
  parent = {MPSI},
  numero = {E},
  niveau = {1},
  titre = {Formation expérimentale},
  chapeau = {
    \cpgeitem{Cette partie est spécifiquement dédiée à la mise en œuvre de la formation expérimentale des étudiants lors des séances de travaux pratiques.}
    \cpgeitem{Dans un premier temps, elle précise les connaissances et savoir-faire qui doivent être acquis dans le domaine de la mesure et de l’évaluation des incertitudes. Elle présente ensuite de façon détaillée l’ensemble des capacités expérimentales qui doivent être acquises en autonomie par les étudiants à l’issue de leur première année de CPGE. Enfin, elle aborde la question de la prévention du risque au laboratoire de physique-chimie.}
    \cpgeitem{Une liste de matériel, que les étudiants doivent savoir utiliser avec l’aide d’une notice succincte, figure dans l’annexe 1 du présent programme.}},
}

\cpgerubrique{MPSI/E/1}{
  parent = {MPSI/E},
  numero = {E.1},
  niveau = {2},
  titre = {Mesures et incertitudes},
  chapeau = {
    \cpgeitem{Les notions et capacités identifiées ci-dessous couvrent les deux années de formation en classe préparatoire aux grandes écoles ; leur pleine maîtrise est donc un objectif de fin de seconde année.}
    \cpgeitem{L’accent est mis sur la variabilité de la mesure d’une grandeur physique et sa caractérisation à l’aide de l’incertitude-type. La comparaison entre deux valeurs mesurées d’une même grandeur physique est conduite au moyen de l’écart normalisé, l’objectif principal étant de développer l’esprit critique des étudiants en s’appuyant sur un critère quantitatif. Le même esprit prévaut dans l’analyse des résultats d’une régression linéaire qui ne saurait s’appuyer sur l’exploitation non raisonnée du coefficient de corrélation ($R^2$).}
    \cpgeitem{Le recours à la simulation vise à illustrer, sur la base de mesures expérimentales, différents effets de la variabilité de la mesure d’une grandeur physique dans les cas des incertitudes-types composées et de la régression linéaire.}},
}

\cpgeligne{MPSI/E/1/1}{
  parent = {MPSI/E/1},
  numero = {E.1.1},
  contenu = {
    \cpgeitem{Variabilité de la mesure d’une grandeur physique. Incertitude.}
    \cpgeitem{Incertitude-type.}},
  capacites = {
    \cpgeitem{Identifier les incertitudes liées, par exemple, à l’opérateur, à l’environnement, aux instruments ou à la méthode de mesure.}
    \cpgeitem{Procéder à l’évaluation d’une incertitude-type par une approche statistique (évaluation de type A).}
    \cpgeitem{Procéder à l’évaluation d’une incertitude-type par une autre approche que statistique (évaluation de type B).}
    \cpgeitem{Associer un intervalle de confiance à l’écart-type dans l’hypothèse d’une distribution suivant la loi normale.}},
  experimentales = {},
  numeriques = {},
}

\cpgeligne{MPSI/E/1/2}{
  parent = {MPSI/E/1},
  numero = {E.1.2},
  contenu = {
    \cpgeitem{Incertitudes-types composées.}},
  capacites = {
    \cpgeitem{Évaluer l’incertitude-type d’une grandeur s’exprimant en fonction d’autres grandeurs, dont les incertitudes-types sont connues, à l’aide d’une somme, d’une différence, d’un produit ou d’un quotient.}
    \cpgeitem{Comparer entre elles les différentes contributions lors de l’évaluation d’une incertitude-type composée.}},
  experimentales = {},
  numeriques = {
    \cpgeitem{simuler, à l’aide d’un langage de programmation ou d’un tableur, un processus aléatoire permettant de caractériser la variabilité de la valeur d’une grandeur composée.}},
}

\cpgeligne{MPSI/E/1/3}{
  parent = {MPSI/E/1},
  numero = {E.1.3},
  contenu = {
    \cpgeitem{Écriture du résultat d’une mesure.}},
  capacites = {
    \cpgeitem{Écrire, avec un nombre adapté de chiffres significatifs, le résultat d’une mesure.}},
  experimentales = {},
  numeriques = {},
}

\cpgeligne{MPSI/E/1/4}{
  parent = {MPSI/E/1},
  numero = {E.1.4},
  contenu = {
    \cpgeitem{Comparaison de deux valeurs ; écart normalisé.}},
  capacites = {
    \cpgeitem{Comparer deux valeurs dont les incertitudes-types sont connues à l’aide de leur écart normalisé.}
    \cpgeitem{Analyser les causes d’une éventuelle incompatibilité entre le résultat d’une mesure et le résultat attendu par une modélisation.}},
  experimentales = {},
  numeriques = {},
}

\cpgeligne{MPSI/E/1/5}{
  parent = {MPSI/E/1},
  numero = {E.1.5},
  contenu = {
    \cpgeitem{Régression linéaire.}},
  capacites = {
    \cpgeitem{Utiliser un logiciel de régression linéaire afin d’obtenir les valeurs des paramètres du modèle.}
    \cpgeitem{Analyser les résultats obtenus à l’aide d’une procédure de validation : analyse graphique intégrant les barres d’incertitude ou analyse des écarts normalisés.}},
  experimentales = {},
  numeriques = {
    \cpgeitem{simuler, à l’aide d’un langage de programmation ou d’un tableur, un processus aléatoire de variation des valeurs expérimentales de l’une des grandeurs – simulation Monte-Carlo – pour évaluer l’incertitude sur les paramètres du modèle.}},
}

\cpgerubrique{MPSI/E/2}{
  parent = {MPSI/E},
  numero = {E.2},
  niveau = {2},
  titre = {Mesures et capacités expérimentales},
  chapeau = {
    \cpgeitem{Cette partie présente l’ensemble des capacités expérimentales que les étudiants doivent acquérir au cours de l’année durant les séances de travaux pratiques. Une séance de travaux pratiques s’articule autour d’une problématique, que les thèmes – repérés en gras dans la colonne « capacités exigibles »   de la partie « Contenus thématiques » du programme – peuvent servir à définir. Le travail de ces capacités et leur consolidation se poursuit en seconde année.}
    \cpgeitem{Dans le tableau ci-dessous, les différentes capacités à acquérir sont groupées par domaines thématiques ou transversaux. Cela ne signifie pas qu’une activité expérimentale se limite à un seul domaine. La capacité à former une image de bonne qualité, par exemple, peut être mobilisée au cours d’une expérience de mécanique ou de thermodynamique, cette transversalité de la formation devant être un moyen, entre d’autres, de favoriser l’autonomie et la prise d’initiative.}},
}

\cpgeligne{MPSI/E/2/1}{
  parent = {MPSI/E/2},
  numero = {E.2.1},
  contenu = {
    \cpgeitem{\textbf{1. Mesures de longueurs et d’angles}}
    \cpgeitem{Longueurs : sur un banc d’optique.}},
  capacites = {
    \cpgeitem{Mettre en œuvre une mesure de longueur par déplacement d’un viseur entre deux positions.}},
  experimentales = {},
  numeriques = {},
}

\cpgeligne{MPSI/E/2/2}{
  parent = {MPSI/E/2},
  numero = {E.2.2},
  contenu = {
    \cpgeitem{Longueurs : à partir d’une photo ou d’une vidéo.}},
  capacites = {
    \cpgeitem{Évaluer, par comparaison à un étalon, une longueur (ou les coordonnées d’une position) sur une image numérique et en estimer la précision.}},
  experimentales = {},
  numeriques = {},
}

\cpgeligne{MPSI/E/2/3}{
  parent = {MPSI/E/2},
  numero = {E.2.3},
  contenu = {
    \cpgeitem{Angles : avec un goniomètre.}},
  capacites = {
    \cpgeitem{Utiliser un viseur à frontale fixe, une lunette autocollimatrice.}
    \cpgeitem{Utiliser des vis micrométriques et un réticule.}},
  experimentales = {},
  numeriques = {},
}

\cpgeligne{MPSI/E/2/4}{
  parent = {MPSI/E/2},
  numero = {E.2.4},
  contenu = {
    \cpgeitem{Longueurs d’onde.}},
  capacites = {
    \cpgeitem{Étudier un spectre à l’aide d’un spectromètre à fibre optique.}
    \cpgeitem{Mesurer une longueur d’onde optique à l’aide d’un goniomètre à réseau.}
    \cpgeitem{Mesurer une longueur d’onde acoustique à l’aide d’un support gradué et d’un oscilloscope bicourbe.}},
  experimentales = {},
  numeriques = {},
}

\cpgeligne{MPSI/E/2/5}{
  parent = {MPSI/E/2},
  numero = {E.2.5},
  contenu = {
    \cpgeitem{\textbf{2. Mesures de temps et de fréquences}}
    \cpgeitem{Fréquence ou période : mesure au fréquencemètre numérique, à l’oscilloscope ou \emph{via} une carte d’acquisition.}},
  capacites = {
    \cpgeitem{Mettre en œuvre une méthode de mesure de fréquence ou de période.}},
  experimentales = {},
  numeriques = {},
}

\cpgeligne{MPSI/E/2/6}{
  parent = {MPSI/E/2},
  numero = {E.2.6},
  contenu = {
    \cpgeitem{Analyse spectrale.}},
  capacites = {
    \cpgeitem{Choisir de façon cohérente la fréquence d’échantillonnage et la durée totale d’acquisition. Effectuer l’analyse spectrale d’un signal périodique à l’aide d’un oscilloscope numérique ou d’une carte d’acquisition.}},
  experimentales = {},
  numeriques = {},
}

\cpgeligne{MPSI/E/2/7}{
  parent = {MPSI/E/2},
  numero = {E.2.7},
  contenu = {
    \cpgeitem{Décalage temporel/déphasage à l’aide d’un oscilloscope numérique.}},
  capacites = {
    \cpgeitem{Reconnaître une avance ou un retard de phase. Passer d’un décalage temporel à un déphasage et inversement.}
    \cpgeitem{Repérer précisément le passage par un déphasage de $0$ ou $\pi$ en mode XY.}},
  experimentales = {},
  numeriques = {},
}

\cpgeligne{MPSI/E/2/8}{
  parent = {MPSI/E/2},
  numero = {E.2.8},
  contenu = {
    \cpgeitem{\textbf{3. Électricité}}
    \cpgeitem{Mesurer une tension : \begin{itemize} \item mesure directe au voltmètre numérique ou à l’oscilloscope numérique. \end{itemize}}
    \cpgeitem{Mesurer l’intensité d’un courant : \begin{itemize} \item mesure directe à l’ampèremètre numérique ; \item mesure indirecte à l’oscilloscope aux bornes d’une résistance adaptée. \end{itemize}}
    \cpgeitem{Mesurer une résistance ou une impédance : \begin{itemize} \item mesure directe à l’ohmmètre/capacimètre ; \item mesure indirecte à l’oscilloscope ou au voltmètre sur un diviseur de tension. \end{itemize}}},
  capacites = {
    \cpgeitem{Capacités communes à l’ensemble des mesures électriques : \begin{itemize} \item expliquer le lien entre résolution, calibre, nombre de points de mesure ; \item préciser la perturbation induite par l’appareil de mesure sur le montage et ses limites (bande passante, résistance d’entrée) ; \item définir la nature de la mesure effectuée (valeur efficace, valeur moyenne, amplitude, valeur crête à crête, etc.). \end{itemize}}},
  experimentales = {},
  numeriques = {},
}

\cpgeligne{MPSI/E/2/9}{
  parent = {MPSI/E/2},
  numero = {E.2.9},
  contenu = {
    \cpgeitem{Produire un signal électrique analogique périodique simple à l’aide d’un GBF.}},
  capacites = {
    \cpgeitem{Obtenir un signal de valeur moyenne, de forme, d’amplitude et de fréquence données.}},
  experimentales = {},
  numeriques = {},
}

\cpgeligne{MPSI/E/2/10}{
  parent = {MPSI/E/2},
  numero = {E.2.10},
  contenu = {
    \cpgeitem{Agir sur un signal électrique à l’aide des fonctions simples suivantes : \begin{itemize} \item isolation, amplification, filtrage ; \item sommation, intégration. \end{itemize}}},
  capacites = {
    \cpgeitem{Gérer, dans un circuit électronique, les contraintes liées à la liaison entre les masses.}
    \cpgeitem{Mettre en œuvre les fonctions de base de l’électronique réalisées par des blocs dont la structure ne fait pas l’objet d’une étude spécifique. Associer ces fonctions de base pour réaliser une fonction complexe en gérant les contraintes liées aux impédances d’entrée et/ou de sortie des blocs.}},
  experimentales = {},
  numeriques = {},
}

\cpgeligne{MPSI/E/2/11}{
  parent = {MPSI/E/2},
  numero = {E.2.11},
  contenu = {
    \cpgeitem{\textbf{4. Optique}}
    \cpgeitem{Former une image.}},
  capacites = {
    \cpgeitem{Éclairer un objet de manière adaptée.}
    \cpgeitem{Choisir une ou plusieurs lentilles en fonction des contraintes expérimentales, et choisir leur focale de façon raisonnée.}
    \cpgeitem{Optimiser la qualité d’une image (alignement, limitation des aberrations, etc.).}
    \cpgeitem{Estimer une valeur approchée d’une distance focale.}},
  experimentales = {},
  numeriques = {},
}

\cpgeligne{MPSI/E/2/12}{
  parent = {MPSI/E/2},
  numero = {E.2.12},
  contenu = {
    \cpgeitem{Créer ou repérer une direction de référence.}},
  capacites = {
    \cpgeitem{Régler et mettre en œuvre une lunette autocollimatrice et un collimateur.}},
  experimentales = {},
  numeriques = {},
}

\cpgeligne{MPSI/E/2/13}{
  parent = {MPSI/E/2},
  numero = {E.2.13},
  contenu = {
    \cpgeitem{Analyser une image numérique.}},
  capacites = {
    \cpgeitem{Acquérir (webcam, appareil photo numérique, etc.) l’image d’un phénomène physique sous forme numérique, et l’exploiter à l’aide d’un logiciel pour conduire l’étude d’un phénomène.}},
  experimentales = {},
  numeriques = {},
}

\cpgeligne{MPSI/E/2/14}{
  parent = {MPSI/E/2},
  numero = {E.2.14},
  contenu = {
    \cpgeitem{\textbf{5. Mécanique}}
    \cpgeitem{Mesurer une masse, un moment d’inertie.}},
  capacites = {
    \cpgeitem{Utiliser une balance de précision.}
    \cpgeitem{Repérer la position d’un centre de masse et mesurer un moment d’inertie à partir d’une période.}},
  experimentales = {},
  numeriques = {},
}

\cpgeligne{MPSI/E/2/15}{
  parent = {MPSI/E/2},
  numero = {E.2.15},
  contenu = {
    \cpgeitem{Visualiser et décomposer un mouvement.}},
  capacites = {
    \cpgeitem{Mettre en œuvre une méthode de stroboscopie. Enregistrer un phénomène à l’aide d’une caméra numérique et repérer la trajectoire à l’aide d’un logiciel dédié, en déduire la vitesse et l’accélération.}},
  experimentales = {},
  numeriques = {},
}

\cpgeligne{MPSI/E/2/16}{
  parent = {MPSI/E/2},
  numero = {E.2.16},
  contenu = {
    \cpgeitem{Mesurer une accélération.}},
  capacites = {
    \cpgeitem{Mettre en œuvre un accéléromètre, par exemple avec l’aide d’un microcontrôleur.}},
  experimentales = {},
  numeriques = {},
}

\cpgeligne{MPSI/E/2/17}{
  parent = {MPSI/E/2},
  numero = {E.2.17},
  contenu = {
    \cpgeitem{Quantifier une action.}},
  capacites = {
    \cpgeitem{Utiliser un dynamomètre.}},
  experimentales = {},
  numeriques = {},
}

\cpgeligne{MPSI/E/2/18}{
  parent = {MPSI/E/2},
  numero = {E.2.18},
  contenu = {
    \cpgeitem{\textbf{6. Thermodynamique}}
    \cpgeitem{Mesurer une pression.}},
  capacites = {
    \cpgeitem{Mettre en œuvre un capteur, en identifiant son caractère différentiel ou absolu.}},
  experimentales = {},
  numeriques = {},
}

\cpgeligne{MPSI/E/2/19}{
  parent = {MPSI/E/2},
  numero = {E.2.19},
  contenu = {
    \cpgeitem{Mesurer une température.}},
  capacites = {
    \cpgeitem{Mettre en œuvre un capteur de température, par exemple avec l’aide d’un microcontrôleur.}
    \cpgeitem{Mettre en œuvre un capteur infrarouge.}
    \cpgeitem{Choisir le capteur en fonction de ses caractéristiques (linéarité, sensibilité, gamme de fonctionnement, temps de réponse), et du type de mesures à effectuer.}},
  experimentales = {},
  numeriques = {},
}

\cpgeligne{MPSI/E/2/20}{
  parent = {MPSI/E/2},
  numero = {E.2.20},
  contenu = {
    \cpgeitem{Effectuer des bilans d’énergie.}},
  capacites = {
    \cpgeitem{Mettre en œuvre une technique de calorimétrie.}},
  experimentales = {},
  numeriques = {},
}

\cpgeligne{MPSI/E/2/21}{
  parent = {MPSI/E/2},
  numero = {E.2.21},
  contenu = {
    \cpgeitem{\textbf{7. Mesures de grandeurs en chimie}}
    \cpgeitem{Mesurer un volume, une masse, un pH, une conductance et une conductivité, une absorbance.}},
  capacites = {
    \cpgeitem{Sélectionner et utiliser le matériel adapté à la précision requise.}
    \cpgeitem{Distinguer les instruments de verrerie In et Ex. Préparer une solution de concentration en masse ou en quantité de matière donnée à partir d’un solide, d’un liquide, d’une solution de composition connue avec le matériel approprié.}
    \cpgeitem{Utiliser les méthodes et le matériel adéquats pour transférer l’intégralité du solide ou du liquide pesé. Utiliser les appareils de mesure (masse, pH, conductance) en s’aidant d’une notice.}
    \cpgeitem{Étalonner une chaîne de mesure si nécessaire.}},
  experimentales = {},
  numeriques = {},
}

\cpgeligne{MPSI/E/2/22}{
  parent = {MPSI/E/2},
  numero = {E.2.22},
  contenu = {
    \cpgeitem{\textbf{8. Analyses qualitatives et quantitatives}}
    \cpgeitem{Effectuer des tests qualitatifs.}},
  capacites = {
    \cpgeitem{Proposer ou mettre en œuvre, à partir d’informations fournies, des tests qualitatifs préalables à l’élaboration d’un protocole.}},
  experimentales = {},
  numeriques = {},
}

\cpgeligne{MPSI/E/2/23}{
  parent = {MPSI/E/2},
  numero = {E.2.23},
  contenu = {
    \cpgeitem{Réaliser des dosages par étalonnage.}},
  capacites = {
    \cpgeitem{Déterminer une concentration en exploitant la mesure de grandeurs physiques caractéristiques de l’espèce ou en construisant et en utilisant une courbe d’étalonnage.}
    \cpgeitem{Déterminer une concentration ou une quantité de matière par spectrophotométrie UV-Visible.}},
  experimentales = {},
  numeriques = {},
}

\cpgeligne{MPSI/E/2/24}{
  parent = {MPSI/E/2},
  numero = {E.2.24},
  contenu = {
    \cpgeitem{Réaliser des dosages par titrage. Titrages directs, indirects. Équivalence. Titrages simples, successifs, simultanés.}
    \cpgeitem{Méthodes expérimentales de suivi d’un titrage : pH-métrie, conductimétrie, indicateurs colorés de fin de titrage.}},
  capacites = {
    \cpgeitem{Identifier et exploiter la réaction support du titrage (recenser les espèces présentes dans le milieu au cours du titrage, repérer l’équivalence, justifier qualitativement l’allure de la courbe ou le changement de couleur observé).}
    \cpgeitem{Proposer ou justifier le protocole d’un titrage à l’aide de données fournies ou à rechercher.}
    \cpgeitem{Mettre en œuvre un protocole expérimental correspondant à un titrage direct ou indirect. Choisir et utiliser un indicateur coloré de fin de titrage.}},
  experimentales = {},
  numeriques = {},
}

\cpgeligne{MPSI/E/2/25}{
  parent = {MPSI/E/2},
  numero = {E.2.25},
  contenu = {
    \cpgeitem{Exploiter des courbes expérimentales de titrage.}},
  capacites = {
    \cpgeitem{Exploiter une courbe de titrage pour déterminer la concentration en espèce titrée.}
    \cpgeitem{Utiliser un logiciel de simulation pour déterminer des courbes de distribution et confronter la courbe de titrage simulée à la courbe expérimentale.}
    \cpgeitem{Distinguer l’équivalence et le repérage du virage d’un indicateur coloré de fin de titrage.}},
  experimentales = {},
  numeriques = {},
}

\cpgeligne{MPSI/E/2/26}{
  parent = {MPSI/E/2},
  numero = {E.2.26},
  contenu = {
    \cpgeitem{Mettre en œuvre des suivis cinétiques de transformations chimiques.}
    \cpgeitem{Suivi en continu de l’évolution temporelle d’une grandeur physique.}},
  capacites = {
    \cpgeitem{Exploiter les résultats d’un suivi temporel de concentration pour déterminer les caractéristiques cinétiques d’une réaction.}
    \cpgeitem{Proposer et mettre en œuvre des conditions expérimentales permettant la simplification de la loi de vitesse.}
    \cpgeitem{Déterminer la valeur d’une énergie d’activation.}},
  experimentales = {},
  numeriques = {},
}

\cpgerubrique{MPSI/E/3}{
  parent = {MPSI/E},
  numero = {E.3},
  niveau = {2},
  titre = {Prévention du risque au laboratoire de physique-chimie},
  chapeau = {
    \cpgeitem{Les étudiants doivent prendre conscience du risque lié à la manipulation et au rejet des produits chimiques. L’apprentissage et le respect des règles de sécurité chimique, électrique, optique et celles liées à la pression et à la température leur permettent de prévenir et de minimiser ce risque.}
    \cpgeitem{Futurs ingénieurs, chercheurs, enseignants, ils doivent être sensibilisés au respect de la législation et à l’impact de leur activité sur l’environnement.}},
}

\cpgeligne{MPSI/E/3/1}{
  parent = {MPSI/E/3},
  numero = {E.3.1},
  contenu = {
    \cpgeitem{\textbf{1. Prévention des risques au laboratoire}}},
  capacites = {
    \cpgeitem{Adopter une attitude responsable et adaptée au travail en laboratoire.}
    \cpgeitem{Développer une attitude autonome dans la prévention des risques.}},
  experimentales = {},
  numeriques = {},
}

\cpgeligne{MPSI/E/3/2}{
  parent = {MPSI/E/3},
  numero = {E.3.2},
  contenu = {
    \cpgeitem{\textbf{Risque chimique}}
    \cpgeitem{Règles de sécurité au laboratoire.}
    \cpgeitem{Classes et catégories de danger.}
    \cpgeitem{Pictogrammes de sécurité pour les produits chimiques.}
    \cpgeitem{Mentions de danger (H) et conseils de prudence (P).}
    \cpgeitem{Fiches de sécurité.}},
  capacites = {
    \cpgeitem{Relever les indications sur le risque associé au prélèvement, au mélange et au stockage des produits chimiques et adopter une attitude responsable lors de leur utilisation.}},
  experimentales = {},
  numeriques = {},
}

\cpgeligne{MPSI/E/3/3}{
  parent = {MPSI/E/3},
  numero = {E.3.3},
  contenu = {
    \cpgeitem{\textbf{Risque électrique}}},
  capacites = {
    \cpgeitem{Adopter une attitude responsable lors de l’utilisation d’appareils électriques.}},
  experimentales = {},
  numeriques = {},
}

\cpgeligne{MPSI/E/3/4}{
  parent = {MPSI/E/3},
  numero = {E.3.4},
  contenu = {
    \cpgeitem{\textbf{Risque optique}}},
  capacites = {
    \cpgeitem{Utiliser les sources laser et les diodes électroluminescentes de manière adaptée.}},
  experimentales = {},
  numeriques = {},
}

\cpgeligne{MPSI/E/3/5}{
  parent = {MPSI/E/3},
  numero = {E.3.5},
  contenu = {
    \cpgeitem{\textbf{Risques liés à la pression et à la température}}},
  capacites = {
    \cpgeitem{Adopter une attitude responsable lors de manipulations de corps chauds ou de dispositifs engageant des hautes ou des}
    \cpgeitem{basses pressions.}},
  experimentales = {},
  numeriques = {},
}

\cpgeligne{MPSI/E/3/6}{
  parent = {MPSI/E/3},
  numero = {E.3.6},
  contenu = {
    % Le texte officiel numérote « 4. » cette rubrique, qui suit la « 1. » :
    % coquille corrigée en « 2. » (voir la documentation, « Écarts au texte officiel »).
    \cpgeitem{\textbf{2. Prévention de l’impact environnemental}}
    \cpgeitem{Traitement et rejet des espèces chimiques.}},
  capacites = {
    \cpgeitem{Adapter le mode d’élimination d’une espèce chimique ou d’un mélange en fonction des informations recueillies sur la toxicité ou les risques.}
    \cpgeitem{Sélectionner, parmi plusieurs modes opératoires, celui qui minimise les impacts environnementaux.}},
  experimentales = {},
  numeriques = {},
}

\cpgerubrique{MPSI/1}{
  parent = {MPSI},
  numero = {1},
  niveau = {1},
  titre = {Ondes et signaux},
  chapeau = {},
}

\cpgerubrique{MPSI/1/1}{
  parent = {MPSI/1},
  numero = {1.1},
  niveau = {2},
  titre = {Formation des images},
  chapeau = {
    \cpgeitem{La partie \cpgecettepartie traite de la formation des images et propose une ouverture sur la notion de guidage de la lumière par une fibre optique. Cette partie est l’occasion d’interroger le concept de modèle en physique et d’en identifier les limites de validité. Elle permet également d’aborder de nombreuses applications technologiques ; certaines sont précisées par le programme, d’autres sont laissées à l’appréciation des enseignants (lunette, microscope, optique d’un smartphone, etc.).}
    \cpgeitem{L’approche expérimentale doit être privilégiée dans ce domaine de la physique qui s’y prête particulièrement bien.}},
}

\cpgeligne{MPSI/1/1/1}{
  parent = {MPSI/1/1},
  numero = {1.1.1},
  contenu = {
    \cpgeitem{\cpgenotion{Sources lumineuses}}
    \cpgeitem{Modèle de la source ponctuelle monochromatique.}
    \cpgeitem{Spectre.}},
  capacites = {
    \cpgeitem{Caractériser une source lumineuse par son spectre. Relier la longueur d’onde dans le vide et la couleur.}},
  experimentales = {},
  numeriques = {},
}

\cpgeligne{MPSI/1/1/2}{
  parent = {MPSI/1/1},
  numero = {1.1.2},
  contenu = {
    \cpgeitem{\cpgenotion{Modèle de l’optique géométrique}}
    \cpgeitem{Modèle de l’optique géométrique. Notion de rayon lumineux. Indice d’un milieu transparent.}},
  capacites = {
    \cpgeitem{Définir le modèle de l’optique géométrique. Indiquer les limites du modèle de l’optique géométrique.}},
  experimentales = {},
  numeriques = {},
}

\cpgeligne{MPSI/1/1/3}{
  parent = {MPSI/1/1},
  numero = {1.1.3},
  contenu = {
    \cpgeitem{Réflexion, réfraction. Lois de Snell-Descartes.}},
  capacites = {
    \cpgeitem{Établir la condition de réflexion totale.}},
  experimentales = {},
  numeriques = {},
}

\cpgeligne{MPSI/1/1/4}{
  parent = {MPSI/1/1},
  numero = {1.1.4},
  contenu = {
    \cpgeitem{\cpgenotion{Conditions de l’approximation de Gauss et applications}}
    \cpgeitem{Stigmatisme. Miroir plan.}},
  capacites = {
    \cpgeitem{Construire l’image d’un objet par un miroir plan.}},
  experimentales = {},
  numeriques = {},
}

\cpgeligne{MPSI/1/1/5}{
  parent = {MPSI/1/1},
  numero = {1.1.5},
  contenu = {
    \cpgeitem{Conditions de l’approximation de Gauss.}},
  capacites = {
    \cpgeitem{Énoncer les conditions de l’approximation de Gauss et ses conséquences.}
    \cpgeitem{Relier le stigmatisme approché aux caractéristiques d’un détecteur.}},
  experimentales = {},
  numeriques = {},
}

\cpgeligne{MPSI/1/1/6}{
  parent = {MPSI/1/1},
  numero = {1.1.6},
  contenu = {
    \cpgeitem{Lentilles minces dans l’approximation de Gauss.}},
  capacites = {
    \cpgeitem{Définir les propriétés du centre optique, des foyers principaux et secondaires, de la distance focale, de la vergence.}
    \cpgeitem{Construire l’image d’un objet situé à distance finie ou infinie à l’aide de rayons lumineux, identifier sa nature réelle ou virtuelle.}
    \cpgeitem{Exploiter les formules de conjugaison et de grandissement transversal de Descartes et de Newton.}
    \cpgeitem{Établir et utiliser la condition de formation de l’image réelle d’un objet réel par une lentille convergente.}},
  experimentales = {},
  numeriques = {},
}

\cpgeligne{MPSI/1/1/7}{
  parent = {MPSI/1/1},
  numero = {1.1.7},
  contenu = {
    \cpgeitem{\cpgenotion{Modèles de quelques dispositifs optiques}}
    \cpgeitem{L’œil.}
    \cpgeitem{Punctum proximum, punctum remotum.}},
  capacites = {
    \cpgeitem{Modéliser l’œil comme l’association d’une lentille de vergence variable et d’un capteur plan fixe.}
    \cpgeitem{Citer les ordres de grandeur de la limite de résolution angulaire et de la plage d’accommodation.}},
  experimentales = {},
  numeriques = {},
}

\cpgeligne{MPSI/1/1/8}{
  parent = {MPSI/1/1},
  numero = {1.1.8},
  contenu = {
    \cpgeitem{L’appareil photographique.}},
  capacites = {
    \cpgeitem{Modéliser l’appareil photographique comme l’association d’une lentille et d’un capteur.}
    \cpgeitem{Construire géométriquement la profondeur de champ pour un réglage donné.}},
  experimentales = {
    \cpgeitem{Étudier l’influence de la focale, de la durée d’exposition, du diaphragme sur la formation de l’image.}},
  numeriques = {},
}

\cpgeligne{MPSI/1/1/9}{
  parent = {MPSI/1/1},
  numero = {1.1.9},
  contenu = {
    \cpgeitem{La fibre optique à saut d’indice.}},
  capacites = {
    \cpgeitem{Établir les expressions du cône d’acceptance et de la dispersion intermodale d’une fibre à saut d’indice.}},
  experimentales = {},
  numeriques = {},
}

\cpgerubrique{MPSI/1/2}{
  parent = {MPSI/1},
  numero = {1.2},
  niveau = {2},
  titre = {Signaux électriques dans l’ARQS},
  chapeau = {
    \cpgeitem{La partie \cpgecettepartie pose les bases nécessaires à l’étude des circuits dans l’Approximation des Régimes Quasi Stationnaires (ARQS). Si le programme se concentre sur l'étude des dipôles $R$, $L$ et $C$, il est possible, lors des travaux pratiques, de faire appel à des composants intégrés ou non linéaires (filtres à capacité commutée, échantillonneur-bloqueur, diodes, photorésistances, etc.) dès lors qu’aucune connaissance préalable n’est nécessaire.}},
}

\cpgeligne{MPSI/1/2/1}{
  parent = {MPSI/1/2},
  numero = {1.2.1},
  contenu = {
    \cpgeitem{Charge électrique, intensité du courant. Potentiel, référence de potentiel, tension. Puissance.}},
  capacites = {
    \cpgeitem{Justifier que l’utilisation de grandeurs électriques continues est compatible avec la quantification de la charge électrique.}
    \cpgeitem{Exprimer l’intensité du courant électrique en termes de débit de charge.}
    \cpgeitem{Exprimer la condition d’application de l’ARQS en fonction de la taille du circuit et de la fréquence.}
    \cpgeitem{Relier la loi des nœuds au postulat de la conservation de la charge.}
    \cpgeitem{Utiliser la loi des mailles.}
    \cpgeitem{Algébriser les grandeurs électriques et utiliser les conventions récepteur et générateur.}
    \cpgeitem{Citer les ordres de grandeur des intensités et des tensions dans différents domaines d’application.}},
  experimentales = {},
  numeriques = {},
}

\cpgeligne{MPSI/1/2/2}{
  parent = {MPSI/1/2},
  numero = {1.2.2},
  contenu = {
    \cpgeitem{Dipôles : résistances, condensateurs, bobines, sources décrites par un modèle linéaire.}},
  capacites = {
    \cpgeitem{Utiliser les relations entre l’intensité et la tension. Citer des ordres de grandeurs des composants $R$, $L$, $C$.}
    \cpgeitem{Exprimer la puissance dissipée par effet Joule dans une résistance.}
    \cpgeitem{Exprimer l’énergie stockée dans un condensateur ou une bobine.}
    \cpgeitem{Modéliser une source en utilisant la représentation de Thévenin.}},
  experimentales = {},
  numeriques = {},
}

\cpgeligne{MPSI/1/2/3}{
  parent = {MPSI/1/2},
  numero = {1.2.3},
  contenu = {
    \cpgeitem{Association de deux résistances.}},
  capacites = {
    \cpgeitem{Remplacer une association série ou parallèle de deux résistances par une résistance équivalente.}
    \cpgeitem{Établir et exploiter les relations des diviseurs de tension ou de courant.}},
  experimentales = {},
  numeriques = {},
}

\cpgeligne{MPSI/1/2/4}{
  parent = {MPSI/1/2},
  numero = {1.2.4},
  contenu = {
    \cpgeitem{Résistance de sortie, résistance d’entrée.}},
  capacites = {},
  experimentales = {
    \cpgeitem{Évaluer une résistance d’entrée ou de sortie à l’aide d’une notice ou d’un appareil afin d’appréhender les conséquences de leurs valeurs sur le fonctionnement d’un circuit.}
    \cpgeitem{Étudier l’influence des résistances d’entrée ou de sortie sur le signal délivré par un GBF, sur la mesure effectuée par un oscilloscope ou un multimètre.}},
  numeriques = {},
}

\cpgerubrique{MPSI/1/3}{
  parent = {MPSI/1},
  numero = {1.3},
  niveau = {2},
  titre = {Circuit linéaire du premier ordre},
  chapeau = {
    \cpgeitem{Les deux parties \cpgecettepartie et \cpgecite{MPSI/1/4} abordent l’étude des circuits linéaires du premier et du second ordre en régime libre puis forcé. Il s’agit avant tout de comprendre les principes des méthodes mises en œuvre et leur exploitation pour étudier le comportement d’un signal traversant un système linéaire. Le choix a été fait de présenter simultanément les oscillateurs électriques et mécaniques de manière à mettre l’accent sur les analogies formelles et comportementales.}},
}

\cpgeligne{MPSI/1/3/1}{
  parent = {MPSI/1/3},
  numero = {1.3.1},
  contenu = {
    \cpgeitem{Régime libre, réponse à un échelon de tension.}},
  capacites = {
    \cpgeitem{Distinguer, sur un relevé expérimental, régime transitoire et régime permanent au cours de l’évolution d’un système du premier ordre soumis à un échelon de tension.}
    \cpgeitem{Interpréter et utiliser la continuité de la tension aux bornes d’un condensateur ou de l’intensité du courant traversant une bobine.}
    \cpgeitem{Établir l’équation différentielle du premier ordre vérifiée par une grandeur électrique dans un circuit comportant une ou deux mailles.}
    \cpgeitem{Déterminer la réponse temporelle dans le cas d’un régime libre ou d’un échelon de tension.}
    \cpgeitem{Déterminer un ordre de grandeur de la durée du régime transitoire.}},
  experimentales = {
    \cpgeitem{Réaliser l’acquisition d’un régime transitoire pour un circuit linéaire du premier ordre et analyser ses caractéristiques. Confronter les résultats expérimentaux aux expressions théoriques.}},
  numeriques = {
    \cpgeitem{mettre en œuvre la méthode d’Euler à l’aide d’un langage de programmation pour simuler la réponse d’un système linéaire du premier ordre à une excitation de forme quelconque.}},
}

\cpgeligne{MPSI/1/3/2}{
  parent = {MPSI/1/3},
  numero = {1.3.2},
  contenu = {
    \cpgeitem{Stockage et dissipation d’énergie.}},
  capacites = {
    \cpgeitem{Réaliser un bilan énergétique.}},
  experimentales = {},
  numeriques = {},
}

\cpgerubrique{MPSI/1/4}{
  parent = {MPSI/1},
  numero = {1.4},
  niveau = {2},
  titre = {Oscillateurs libres et forcés},
  chapeau = {
    \cpgeitem{\cpgerenvoi{MPSI/1/4}}},
}

\cpgeligne{MPSI/1/4/1}{
  parent = {MPSI/1/4},
  numero = {1.4.1},
  contenu = {
    \cpgeitem{Oscillateur harmonique. Exemples du circuit LC et de l’oscillateur mécanique.}},
  capacites = {
    \cpgeitem{Établir et reconnaître l’équation différentielle qui caractérise un oscillateur harmonique ; la résoudre compte tenu des conditions initiales.}
    \cpgeitem{Caractériser l’évolution en utilisant les notions d’amplitude, de phase, de période, de fréquence, de pulsation.}
    \cpgeitem{Réaliser un bilan énergétique.}},
  experimentales = {},
  numeriques = {},
}

\cpgeligne{MPSI/1/4/2}{
  parent = {MPSI/1/4},
  numero = {1.4.2},
  contenu = {
    \cpgeitem{Circuit RLC série et oscillateur mécanique amorti par frottement visqueux.}},
  capacites = {
    \cpgeitem{Analyser, sur des relevés expérimentaux, l’évolution de la forme des régimes transitoires en fonction des paramètres caractéristiques.}
    \cpgeitem{Prévoir l’évolution du système à partir de considérations énergétiques.}
    \cpgeitem{Écrire sous forme canonique l’équation différentielle afin d’identifier la pulsation propre et le facteur de qualité.}
    \cpgeitem{Décrire la nature de la réponse en fonction de la valeur du facteur de qualité.}
    \cpgeitem{Déterminer la réponse détaillée dans le cas d’un régime libre ou d’un système soumis à un échelon en recherchant les racines du polynôme caractéristique. Déterminer un ordre de grandeur de la durée du régime transitoire selon la valeur du facteur de qualité.}},
  experimentales = {
    \cpgeitem{Mettre en évidence la similitude des comportements des oscillateurs mécanique et électronique.}
    \cpgeitem{Réaliser l’acquisition d’un régime transitoire pour un système linéaire du deuxième ordre et analyser ses caractéristiques.}},
  numeriques = {},
}

\cpgeligne{MPSI/1/4/3}{
  parent = {MPSI/1/4},
  numero = {1.4.3},
  contenu = {
    \cpgeitem{Stockage et dissipation d’énergie.}},
  capacites = {
    \cpgeitem{Réaliser un bilan énergétique.}},
  experimentales = {},
  numeriques = {},
}

\cpgeligne{MPSI/1/4/4}{
  parent = {MPSI/1/4},
  numero = {1.4.4},
  contenu = {
    \cpgeitem{Impédances complexes.}},
  capacites = {
    \cpgeitem{Établir et connaître l’impédance d’une résistance, d’un condensateur, d’une bobine.}},
  experimentales = {},
  numeriques = {},
}

\cpgeligne{MPSI/1/4/5}{
  parent = {MPSI/1/4},
  numero = {1.4.5},
  contenu = {
    \cpgeitem{Association de deux impédances.}},
  capacites = {
    \cpgeitem{Remplacer une association série ou parallèle de deux impédances par une impédance équivalente.}},
  experimentales = {},
  numeriques = {},
}

\cpgeligne{MPSI/1/4/6}{
  parent = {MPSI/1/4},
  numero = {1.4.6},
  contenu = {
    \cpgeitem{Oscillateur électrique ou mécanique soumis à une excitation sinusoïdale. Résonance.}},
  capacites = {
    \cpgeitem{Utiliser la représentation complexe pour étudier le régime forcé.}
    \cpgeitem{Relier l’acuité d’une résonance au facteur de qualité. Déterminer la pulsation propre et le facteur de qualité à partir de graphes expérimentaux d’amplitude et de phase.}},
  experimentales = {
    \cpgeitem{Mettre en œuvre un dispositif expérimental visant à caractériser un phénomène de résonance.}},
  numeriques = {},
}

\cpgerubrique{MPSI/1/5}{
  parent = {MPSI/1},
  numero = {1.5},
  niveau = {2},
  titre = {Filtrage linéaire},
  chapeau = {
    \cpgeitem{L’objectif principal de la partie \cpgecettepartie n’est pas de former les étudiants aux aspects techniques des calculs des fonctions de transfert et des tracés de diagrammes de Bode mais de mettre l’accent sur l’interprétation des propriétés du signal de sortie connaissant celles du signal d’entrée et d’appréhender le rôle central de la linéarité des systèmes utilisés.}},
}

\cpgeligne{MPSI/1/5/1}{
  parent = {MPSI/1/5},
  numero = {1.5.1},
  contenu = {
    \cpgeitem{Signaux périodiques.}},
  capacites = {
    \cpgeitem{Analyser la décomposition fournie d’un signal périodique en une somme de fonctions sinusoïdales. Définir la valeur moyenne et la valeur efficace d’un signal.}
    \cpgeitem{Établir par le calcul la valeur efficace d’un signal sinusoïdal.}
    \cpgeitem{Interpréter le fait que le carré de la valeur efficace d’un signal périodique est égal à la somme des carrés des valeurs efficaces de ses harmoniques.}},
  experimentales = {},
  numeriques = {},
}

\cpgeligne{MPSI/1/5/2}{
  parent = {MPSI/1/5},
  numero = {1.5.2},
  contenu = {
    \cpgeitem{Fonction de transfert harmonique. Diagramme de Bode.}},
  capacites = {
    \cpgeitem{Tracer le diagramme de Bode (amplitude et phase) associé à une fonction de transfert d’ordre 1.}
    \cpgeitem{Utiliser une fonction de transfert donnée d’ordre 1 ou 2 (ou ses représentations graphiques) pour étudier la réponse d’un système linéaire à une excitation sinusoïdale, à une somme finie d’excitations sinusoïdales, à un signal périodique.}
    \cpgeitem{Utiliser les échelles logarithmiques et interpréter les zones rectilignes des diagrammes de Bode en amplitude d’après l’expression de la fonction de transfert.}},
  experimentales = {
    \cpgeitem{Mettre en œuvre un dispositif expérimental illustrant l’utilité des fonctions de transfert pour un système linéaire à un ou plusieurs étages.}},
  numeriques = {},
}

\cpgeligne{MPSI/1/5/3}{
  parent = {MPSI/1/5},
  numero = {1.5.3},
  contenu = {
    \cpgeitem{Modèles de filtres passifs : passe-bas et passe-haut d’ordre 1, passe-bas et passe-bande d’ordre 2.}},
  capacites = {
    \cpgeitem{Choisir un modèle de filtre en fonction d’un cahier des charges.}
    \cpgeitem{Expliciter les conditions d’utilisation d’un filtre en tant que moyenneur, intégrateur, ou dérivateur.}
    \cpgeitem{Expliquer l’intérêt, pour garantir leur fonctionnement lors de mises en cascade, de réaliser des filtres de tension de faible impédance de sortie et forte impédance d’entrée.}
    \cpgeitem{Expliquer la nature du filtrage introduit par un dispositif mécanique (sismomètre, amortisseur, accéléromètre, etc.).}},
  experimentales = {
    \cpgeitem{Étudier le filtrage linéaire d’un signal non sinusoïdal à partir d’une analyse spectrale.}
    \cpgeitem{Détecter le caractère non linéaire d’un système par l’apparition de nouvelles fréquences.}},
  numeriques = {
    \cpgeitem{simuler, à l’aide d’un langage de programmation, l'action d'un filtre sur un signal périodique dont le spectre est fourni. Mettre en évidence l'influence des caractéristiques du filtre sur l'opération de filtrage.}},
}

\cpgerubrique{MPSI/1/6}{
  parent = {MPSI/1},
  numero = {1.6},
  niveau = {2},
  titre = {Propagation d’un signal},
  chapeau = {
    \cpgeitem{Dans la partie \cpgecettepartie, il est recommandé de s’appuyer sur une approche expérimentale ou sur des logiciels de simulation pour permettre aux étudiants de faire le lien entre l’observation de signaux qui se propagent et la traduction mathématique de cette propagation, sans qu’aucune référence ne soit faite à une équation d’onde. L'étude de la somme de deux signaux sinusoïdaux de même fréquence et du phénomène d’interférences associé permet de mettre en évidence le rôle essentiel joué par le déphasage entre les deux signaux dans le signal résultant. L’étude des interférences lumineuses est l’occasion d’introduire la notion de différence de chemin optique et de la relier au déphasage.}},
}

\cpgeligne{MPSI/1/6/1}{
  parent = {MPSI/1/6},
  numero = {1.6.1},
  contenu = {
    \cpgeitem{Exemples de signaux. Signal sinusoïdal.}},
  capacites = {
    \cpgeitem{Identifier les grandeurs physiques correspondant à des signaux acoustiques, électriques, électromagnétiques.}},
  experimentales = {},
  numeriques = {},
}

\cpgeligne{MPSI/1/6/2}{
  parent = {MPSI/1/6},
  numero = {1.6.2},
  contenu = {
    \cpgeitem{\cpgenotion{Propagation d’un signal dans un milieu illimité, non dispersif et transparent}}
    \cpgeitem{Onde progressive dans le cas d’une propagation unidimensionnelle non dispersive}
    \cpgeitem{Célérité, retard temporel.}},
  capacites = {
    \cpgeitem{Écrire les signaux sous la forme $f(x-ct)$ ou $g(x+ct)$. Écrire les signaux sous la forme $f(t-x/c)$ ou $g(t+x/c)$.}
    \cpgeitem{Prévoir, dans le cas d’une onde progressive, l’évolution temporelle à position fixée et l’évolution spatiale à différents instants.}},
  experimentales = {},
  numeriques = {},
}

\cpgeligne{MPSI/1/6/3}{
  parent = {MPSI/1/6},
  numero = {1.6.3},
  contenu = {
    \cpgeitem{Modèle de l’onde progressive sinusoïdale unidimensionnelle. Vitesse de phase, déphasage, double périodicité spatiale et temporelle.}},
  capacites = {
    \cpgeitem{Citer quelques ordres de grandeur de fréquences dans les domaines acoustique, mécanique et électromagnétique.}
    \cpgeitem{Établir la relation entre la fréquence, la longueur d’onde et la vitesse de phase.}
    \cpgeitem{Relier le déphasage entre les signaux perçus en deux points distincts au retard dû à la propagation.}},
  experimentales = {
    \cpgeitem{Mesurer la vitesse de phase, la longueur d’onde et le déphasage dû à la propagation d’un phénomène ondulatoire.}},
  numeriques = {},
}

\cpgeligne{MPSI/1/6/4}{
  parent = {MPSI/1/6},
  numero = {1.6.4},
  contenu = {
    \cpgeitem{Milieux dispersifs ou non dispersifs.}},
  capacites = {
    \cpgeitem{Définir un milieu dispersif.}
    \cpgeitem{Citer des exemples de situations de propagation dispersive et non dispersive.}},
  experimentales = {},
  numeriques = {},
}

\cpgeligne{MPSI/1/6/5}{
  parent = {MPSI/1/6},
  numero = {1.6.5},
  contenu = {
    \cpgeitem{\cpgenotion{Phénomène d’interférences}}
    \cpgeitem{Interférences entre deux ondes acoustiques ou mécaniques de même fréquence.}},
  capacites = {
    \cpgeitem{Exprimer les conditions d’interférences constructives ou destructives.}
    \cpgeitem{Déterminer l’amplitude de l’onde résultante en un point en fonction du déphasage.}},
  experimentales = {},
  numeriques = {},
}

\cpgeligne{MPSI/1/6/6}{
  parent = {MPSI/1/6},
  numero = {1.6.6},
  contenu = {
    \cpgeitem{Interférences entre deux ondes lumineuses de même fréquence.}
    \cpgeitem{Exemple du dispositif des trous d’Young éclairé par une source monochromatique.}
    \cpgeitem{Différence de chemin optique. Conditions d’interférences constructives ou destructives.}
    \cpgeitem{Formule de Fresnel.}},
  capacites = {
    \cpgeitem{Relier le déphasage entre les deux ondes à la différence de chemin optique.}
    \cpgeitem{Établir l’expression littérale de la différence de chemin optique entre les deux ondes.}
    \cpgeitem{Exploiter la formule de Fresnel fournie pour décrire la répartition d’intensité lumineuse.}},
  experimentales = {
    \cpgeitem{Mettre en œuvre un dispositif expérimental pour visualiser et caractériser le phénomène d’interférences de deux ondes.}},
  numeriques = {},
}

\cpgerubrique{MPSI/1/7}{
  parent = {MPSI/1},
  numero = {1.7},
  niveau = {2},
  titre = {Induction et forces de Laplace},
  chapeau = {
    \cpgeitem{La partie \cpgecettepartie s’appuie sur les nombreuses applications présentes dans notre environnement immédiat : boussole, moteur électrique, alternateur, transformateur, haut-parleur, plaques à induction, carte RFID… Il s’agit de restituer toute la richesse de ces applications dans un volume horaire modeste, ce qui limite les géométries envisagées et le formalisme utilisé. Le point de vue adopté cherche à mettre l’accent sur les phénomènes et sur la modélisation sommaire de leurs applications. Toute étude du champ électromoteur est exclue. L’induction et les forces de Laplace dans un circuit mobile sont introduites dans le cas d’un champ uniforme et stationnaire, soit dans le modèle des rails de Laplace, soit dans celui d’un cadre rectangulaire en rotation. Ce dernier modèle permet d’introduire la notion de dipôle magnétique et une analogie de comportement permet de l’étendre au cas de l’aiguille d’une boussole.}
    \cpgeitem{Le succès de cet enseignement suppose le respect de ces limitations : il ne s’agit pas d’une étude générale des phénomènes d’induction. Corrélativement, l’enseignement de cette partie doit impérativement s’appuyer sur une démarche expérimentale authentique, qu’il s’agisse d’expériences de cours ou d’activités expérimentales.}},
}

\cpgerubrique{MPSI/1/7/1}{
  parent = {MPSI/1/7},
  numero = {1.7.1},
  niveau = {3},
  titre = {Champ magnétique},
  chapeau = {
    \cpgeitem{La partie \cpgecettepartie vise à relier le champ magnétique et ses sources ; l’accent est mis sur le concept de champ vectoriel, l’analyse des symétries et des invariances, l’exploitation des représentations graphiques et la connaissance d’ordres de grandeur.}},
}

\cpgeligne{MPSI/1/7/1/1}{
  parent = {MPSI/1/7/1},
  numero = {1.7.1.1},
  contenu = {
    \cpgeitem{Sources de champ magnétique ; cartes de champ magnétique.}},
  capacites = {
    \cpgeitem{Exploiter une représentation graphique d’un champ vectoriel, identifier les zones de champ uniforme, de champ faible et l’emplacement des sources.}
    \cpgeitem{Tracer l’allure des cartes de champs magnétiques pour un aimant droit, une spire circulaire et une bobine longue.}
    \cpgeitem{Décrire un dispositif permettant de réaliser un champ magnétique quasi uniforme.}
    \cpgeitem{Citer des ordres de grandeur de champs magnétiques : au voisinage d’aimants, dans un appareil d’IRM, dans le cas du champ magnétique terrestre.}},
  experimentales = {},
  numeriques = {},
}

\cpgeligne{MPSI/1/7/1/2}{
  parent = {MPSI/1/7/1},
  numero = {1.7.1.2},
  contenu = {
    \cpgeitem{Symétries et invariances des distributions de courant.}},
  capacites = {
    \cpgeitem{Exploiter les propriétés de symétrie et d’invariance des sources pour prévoir des propriétés du champ créé.}},
  experimentales = {},
  numeriques = {},
}

\cpgeligne{MPSI/1/7/1/3}{
  parent = {MPSI/1/7/1},
  numero = {1.7.1.3},
  contenu = {
    \cpgeitem{Lien entre le champ magnétique et l’intensité du courant.}},
  capacites = {
    \cpgeitem{Évaluer l’ordre de grandeur d’un champ magnétique à partir d’expressions fournies.}},
  experimentales = {},
  numeriques = {},
}

\cpgeligne{MPSI/1/7/1/4}{
  parent = {MPSI/1/7/1},
  numero = {1.7.1.4},
  contenu = {
    \cpgeitem{Moment magnétique.}},
  capacites = {
    \cpgeitem{Définir le moment magnétique associé à une boucle de courant plane.}
    \cpgeitem{Associer à un aimant un moment magnétique par analogie avec une boucle de courant.}
    \cpgeitem{Citer un ordre de grandeur du moment magnétique associé à un aimant usuel.}},
  experimentales = {},
  numeriques = {},
}

\cpgerubrique{MPSI/1/7/2}{
  parent = {MPSI/1/7},
  numero = {1.7.2},
  niveau = {3},
  titre = {Actions d’un champ magnétique},
  chapeau = {
    \cpgeitem{Dans la partie \cpgecettepartie, l’enseignant est libre d’introduire la force de Laplace avec ou sans référence à la force de Lorentz. Il s’agit ici de se doter d’expressions opérationnelles pour étudier le mouvement dans un champ uniforme et stationnaire (soit d’une barre en translation, soit d’un moment magnétique en rotation modélisé par un cadre rectangulaire).}},
}

\cpgeligne{MPSI/1/7/2/1}{
  parent = {MPSI/1/7/2},
  numero = {1.7.2.1},
  contenu = {
    \cpgeitem{Densité linéique de la force de Laplace dans le cas d’un élément de courant filiforme.}},
  capacites = {
    \cpgeitem{Différencier le champ magnétique extérieur subi du champ magnétique propre créé par le courant filiforme.}},
  experimentales = {},
  numeriques = {},
}

\cpgeligne{MPSI/1/7/2/2}{
  parent = {MPSI/1/7/2},
  numero = {1.7.2.2},
  contenu = {
    \cpgeitem{Résultante et puissance des forces de Laplace.}},
  capacites = {
    \cpgeitem{Établir et citer l’expression de la résultante des forces de Laplace dans le cas d’une barre conductrice placée dans un champ magnétique extérieur uniforme et stationnaire.}
    \cpgeitem{Exprimer la puissance des forces de Laplace.}},
  experimentales = {},
  numeriques = {},
}

\cpgeligne{MPSI/1/7/2/3}{
  parent = {MPSI/1/7/2},
  numero = {1.7.2.3},
  contenu = {
    \cpgeitem{Couple et puissance des actions mécaniques de Laplace dans le cas d’une spire rectangulaire, parcourue par un courant, en rotation autour d’un axe de symétrie de la spire passant par les deux milieux de côtés opposés et placée dans un champ magnétique extérieur uniforme et stationnaire orthogonal à l’axe.}},
  capacites = {
    \cpgeitem{Établir et exploiter l’expression du moment du couple subi en fonction du champ magnétique extérieur et du moment magnétique.}
    \cpgeitem{Exprimer la puissance des actions mécaniques de Laplace.}},
  experimentales = {},
  numeriques = {},
}

\cpgeligne{MPSI/1/7/2/4}{
  parent = {MPSI/1/7/2},
  numero = {1.7.2.4},
  contenu = {
    \cpgeitem{Action d’un champ magnétique extérieur uniforme sur un aimant.}
    \cpgeitem{Positions d’équilibre et stabilité.}},
  capacites = {},
  experimentales = {
    \cpgeitem{Mettre en œuvre un dispositif expérimental pour étudier l’action d’un champ magnétique uniforme sur une boussole.}},
  numeriques = {},
}

\cpgeligne{MPSI/1/7/2/5}{
  parent = {MPSI/1/7/2},
  numero = {1.7.2.5},
  contenu = {
    \cpgeitem{Effet moteur d’un champ magnétique tournant.}},
  capacites = {},
  experimentales = {
    \cpgeitem{Créer un champ magnétique tournant à l’aide de deux ou trois bobines et mettre en rotation une aiguille aimantée.}},
  numeriques = {},
}

\cpgerubrique{MPSI/1/7/3}{
  parent = {MPSI/1/7},
  numero = {1.7.3},
  niveau = {3},
  titre = {Lois de l’induction},
  chapeau = {
    \cpgeitem{La partie \cpgecettepartie repose sur la loi de Faraday qui se prête parfaitement à une introduction expérimentale et qui constitue un bel exemple d’illustration de l’histoire des sciences. On évoque, à ce sujet, les différents points de vue possibles sur le même phénomène selon le référentiel dans lequel on se place.}},
}

\cpgeligne{MPSI/1/7/3/1}{
  parent = {MPSI/1/7/3},
  numero = {1.7.3.1},
  contenu = {
    \cpgeitem{\cpgenotion{Flux d’un champ magnétique}}
    \cpgeitem{Flux d’un champ magnétique à travers une surface s’appuyant sur un contour fermé orienté.}},
  capacites = {
    \cpgeitem{Évaluer le flux d’un champ magnétique uniforme à}
    \cpgeitem{travers une surface s’appuyant sur un contour fermé orienté plan.}},
  experimentales = {},
  numeriques = {},
}

\cpgeligne{MPSI/1/7/3/2}{
  parent = {MPSI/1/7/3},
  numero = {1.7.3.2},
  contenu = {
    \cpgeitem{\cpgenotion{Loi de Faraday}}
    \cpgeitem{Courant induit par le déplacement relatif d’une boucle conductrice par rapport à un aimant ou un circuit inducteur. Sens du courant induit.}},
  capacites = {},
  experimentales = {
    \cpgeitem{Décrire, mettre en œuvre et interpréter des expériences illustrant les lois de Lenz et de Faraday.}},
  numeriques = {},
}

\cpgeligne{MPSI/1/7/3/3}{
  parent = {MPSI/1/7/3},
  numero = {1.7.3.3},
  contenu = {
    \cpgeitem{Loi de modération de Lenz.}},
  capacites = {
    \cpgeitem{Utiliser la loi de Lenz pour prédire ou interpréter les phénomènes physiques observés.}},
  experimentales = {},
  numeriques = {},
}

\cpgeligne{MPSI/1/7/3/4}{
  parent = {MPSI/1/7/3},
  numero = {1.7.3.4},
  contenu = {
    \cpgeitem{Force électromotrice induite, loi de Faraday.}},
  capacites = {
    \cpgeitem{Utiliser la loi de Faraday en précisant les conventions d’algébrisation.}},
  experimentales = {},
  numeriques = {},
}

\cpgerubrique{MPSI/1/7/4}{
  parent = {MPSI/1/7},
  numero = {1.7.4},
  niveau = {3},
  titre = {Circuit fixe dans un champ magnétique qui dépend du temps},
  chapeau = {
    \cpgeitem{La partie \cpgecettepartie aborde le phénomène d’auto-induction puis le couplage par mutuelle inductance entre deux circuits fixes. Elle traite du modèle du transformateur de tensions.}},
}

\cpgeligne{MPSI/1/7/4/1}{
  parent = {MPSI/1/7/4},
  numero = {1.7.4.1},
  contenu = {
    \cpgeitem{\cpgenotion{Auto-induction}}
    \cpgeitem{Flux propre et inductance propre.}},
  capacites = {
    \cpgeitem{Différencier le flux propre des flux extérieurs. Utiliser la loi de modération de Lenz.}
    \cpgeitem{Évaluer et citer l’ordre de grandeur de l’inductance propre d’une bobine de grande longueur.}},
  experimentales = {
    \cpgeitem{Mesurer la valeur de l’inductance propre d’une bobine.}},
  numeriques = {},
}

\cpgeligne{MPSI/1/7/4/2}{
  parent = {MPSI/1/7/4},
  numero = {1.7.4.2},
  contenu = {
    \cpgeitem{Étude énergétique.}},
  capacites = {
    \cpgeitem{Réaliser un bilan de puissance et d’énergie dans un système siège d’un phénomène d’auto-induction en s’appuyant sur un schéma électrique équivalent.}},
  experimentales = {},
  numeriques = {},
}

\cpgeligne{MPSI/1/7/4/3}{
  parent = {MPSI/1/7/4},
  numero = {1.7.4.3},
  contenu = {
    \cpgeitem{\cpgenotion{Cas de deux bobines en interaction}}
    \cpgeitem{Inductance mutuelle entre deux bobines.}},
  capacites = {
    \cpgeitem{Déterminer l’inductance mutuelle entre deux bobines de même axe de grande longueur en «~influence totale~»}},
  experimentales = {},
  numeriques = {},
}

\cpgeligne{MPSI/1/7/4/4}{
  parent = {MPSI/1/7/4},
  numero = {1.7.4.4},
  contenu = {
    \cpgeitem{Circuits électriques à une maille couplés par le phénomène de mutuelle induction en régime sinusoïdal forcé.}},
  capacites = {
    \cpgeitem{Citer des applications dans le domaine de l’industrie ou de la vie courante.}
    \cpgeitem{Établir le système d’équations en régime sinusoïdal forcé en s’appuyant sur des schémas électriques équivalents.}},
  experimentales = {},
  numeriques = {},
}

\cpgeligne{MPSI/1/7/4/5}{
  parent = {MPSI/1/7/4},
  numero = {1.7.4.5},
  contenu = {
    \cpgeitem{Étude énergétique.}},
  capacites = {
    \cpgeitem{Réaliser un bilan de puissance et d’énergie.}},
  experimentales = {},
  numeriques = {},
}

\cpgerubrique{MPSI/1/7/5}{
  parent = {MPSI/1/7},
  numero = {1.7.5},
  niveau = {3},
  titre = {Circuit mobile dans un champ magnétique stationnaire},
  chapeau = {
    \cpgeitem{La partie \cpgecettepartie est centrée sur la conversion de puissance. Des situations géométriques simples permettent de dégager les paramètres physiques pertinents afin de modéliser, par exemple, un dispositif de freinage.}},
}

\cpgeligne{MPSI/1/7/5/1}{
  parent = {MPSI/1/7/5},
  numero = {1.7.5.1},
  contenu = {
    \cpgeitem{\cpgenotion{Conversion de puissance mécanique en puissance électrique}}
    \cpgeitem{Rail de Laplace.}
    \cpgeitem{Spire rectangulaire soumise à un champ magnétique extérieur uniforme et en rotation uniforme autour d’un axe fixe orthogonal au champ magnétique.}},
  capacites = {
    \cpgeitem{Interpréter qualitativement les phénomènes observés.}
    \cpgeitem{Écrire les équations électrique et mécanique en précisant les conventions de signe.}
    \cpgeitem{Effectuer un bilan énergétique.}
    \cpgeitem{Citer des applications dans le domaine de l’industrie ou de la vie courante.}},
  experimentales = {},
  numeriques = {},
}

\cpgeligne{MPSI/1/7/5/2}{
  parent = {MPSI/1/7/5},
  numero = {1.7.5.2},
  contenu = {
    \cpgeitem{Freinage par induction.}},
  capacites = {
    \cpgeitem{Expliquer l’origine des courants de Foucault et en citer des exemples d’utilisation.}},
  experimentales = {
    \cpgeitem{Mettre en évidence qualitativement les courants de Foucault.}},
  numeriques = {},
}

\cpgerubrique{MPSI/1/8}{
  parent = {MPSI/1},
  numero = {1.8},
  niveau = {2},
  titre = {Introduction à la physique quantique},
  chapeau = {
    \cpgeitem{La partie \cpgecettepartie est structurée autour de la présentation d’expériences réalisées depuis le début du XXème siècle. Cette partie vise à questionner la représentation classique du monde proposée dans les autres parties du programme. Les concepts essentiels abordés sont la dualité onde-particule, l’interprétation probabiliste de la fonction d’onde, l’inégalité de Heisenberg spatiale et la quantification de l’énergie dans les atomes. La réflexion sur les thèmes abordés ici est avant tout qualitative ; toute dérive calculatoire exploitant les concepts propres à la physique quantique doit être évitée.}},
}

\cpgeligne{MPSI/1/8/1}{
  parent = {MPSI/1/8},
  numero = {1.8.1},
  contenu = {
    \cpgeitem{\cpgenotion{Dualité onde-particule pour la lumière et la matière}}
    \cpgeitem{Photon : énergie et impulsion.}},
  capacites = {
    \cpgeitem{Décrire un exemple d’expérience mettant en évidence la nécessité de la notion de photon.}},
  experimentales = {},
  numeriques = {},
}

\cpgeligne{MPSI/1/8/2}{
  parent = {MPSI/1/8},
  numero = {1.8.2},
  contenu = {
    \cpgeitem{Onde de matière associée à une particule. Relation de de Broglie.}},
  capacites = {
    \cpgeitem{Décrire un exemple d’expérience mettant en évidence le comportement ondulatoire de la matière.}
    \cpgeitem{Évaluer des ordres de grandeurs typiques intervenant dans des phénomènes quantiques.}},
  experimentales = {},
  numeriques = {},
}

\cpgeligne{MPSI/1/8/3}{
  parent = {MPSI/1/8},
  numero = {1.8.3},
  contenu = {
    \cpgeitem{\cpgenotion{Introduction au formalisme quantique}}
    \cpgeitem{Fonction d’onde : introduction qualitative, interprétation probabiliste.}},
  capacites = {
    \cpgeitem{Interpréter une expérience d’interférences (matière ou lumière) «~particule par particule~» en termes probabilistes.}},
  experimentales = {},
  numeriques = {},
}

\cpgeligne{MPSI/1/8/4}{
  parent = {MPSI/1/8},
  numero = {1.8.4},
  contenu = {
    \cpgeitem{Inégalité de Heisenberg spatiale.}},
  capacites = {
    \cpgeitem{Établir par analogie avec la diffraction des ondes lumineuses, l’inégalité en ordre de grandeur :}
    \cpgeitem{$\Delta p\Delta x\geq\hbar$.}},
  experimentales = {},
  numeriques = {},
}

\cpgeligne{MPSI/1/8/5}{
  parent = {MPSI/1/8},
  numero = {1.8.5},
  contenu = {
    \cpgeitem{\cpgenotion{Quantification de l’énergie}}
    \cpgeitem{Modèle planétaire de Bohr. Limites.}},
  capacites = {
    \cpgeitem{Exploiter l’hypothèse de quantification du moment cinétique orbital pour obtenir l’expression des niveaux d’énergie électronique de l’atome d’hydrogène.}},
  experimentales = {},
  numeriques = {},
}

\cpgerubrique{MPSI/2}{
  parent = {MPSI},
  numero = {2},
  niveau = {1},
  titre = {Mouvements et interactions},
  chapeau = {},
}

\cpgerubrique{MPSI/2/1}{
  parent = {MPSI/2},
  numero = {2.1},
  niveau = {2},
  titre = {Description et paramétrage du mouvement d’un point},
  chapeau = {
    \cpgeitem{La partie \cpgecettepartie vise notamment à mettre en place les principaux systèmes de coordonnées : cartésiennes, polaires, cylindriques et sphériques. Le but est de permettre aux étudiants de disposer d’outils efficaces pour décrire une grande variété de mouvements de points. Pour atteindre cet objectif, il convient de les familiariser progressivement avec les projections et dérivations de vecteurs ainsi qu’avec l’algébrisation des grandeurs dans un contexte   relevant de la physique. Enfin, cette partie est l’occasion de procéder à des analyses qualitatives des comportements cinématiques de systèmes réels assimilés à un point, notamment sur les exemples simples des mouvements rectilignes et circulaires.}},
}

\cpgeligne{MPSI/2/1/1}{
  parent = {MPSI/2/1},
  numero = {2.1.1},
  contenu = {
    \cpgeitem{\cpgenotion{Repérage dans l’espace et dans le temps}}
    \cpgeitem{Espace et temps classiques. Notion de référentiel. Caractère relatif du mouvement. Caractère absolu des distances et des intervalles de temps.}},
  capacites = {
    \cpgeitem{Citer une situation où la description classique de l’espace ou du temps est prise en défaut.}},
  experimentales = {},
  numeriques = {},
}

\cpgeligne{MPSI/2/1/2}{
  parent = {MPSI/2/1},
  numero = {2.1.2},
  contenu = {
    \cpgeitem{\cpgenotion{Cinématique du point}}
    \cpgeitem{Description du mouvement d’un point. Vecteurs position, vitesse et accélération.}
    \cpgeitem{Systèmes de coordonnées cartésiennes, cylindriques et sphériques.}},
  capacites = {
    \cpgeitem{Exprimer à partir d’un schéma le déplacement élémentaire dans les différents systèmes de coordonnées, construire le trièdre local associé et en déduire géométriquement les composantes du vecteur vitesse en coordonnées cartésiennes et cylindriques.}
    \cpgeitem{Établir les expressions des composantes des vecteurs position, déplacement élémentaire, vitesse et accélération dans les seuls cas des coordonnées cartésiennes et cylindriques.}
    \cpgeitem{Identifier les degrés de liberté d’un mouvement. Choisir un système de coordonnées adapté au problème.}},
  experimentales = {},
  numeriques = {},
}

\cpgeligne{MPSI/2/1/3}{
  parent = {MPSI/2/1},
  numero = {2.1.3},
  contenu = {
    \cpgeitem{Mouvement à vecteur accélération constant.}},
  capacites = {
    \cpgeitem{Exprimer le vecteur vitesse et le vecteur position en fonction du temps.}
    \cpgeitem{Établir l’expression de la trajectoire en coordonnées}
    \cpgeitem{cartésiennes.}},
  experimentales = {},
  numeriques = {},
}

\cpgeligne{MPSI/2/1/4}{
  parent = {MPSI/2/1},
  numero = {2.1.4},
  contenu = {
    \cpgeitem{Mouvement circulaire uniforme et non uniforme.}},
  capacites = {
    \cpgeitem{Exprimer les composantes du vecteur position, du vecteur vitesse et du vecteur accélération en coordonnées polaires planes.}},
  experimentales = {},
  numeriques = {},
}

\cpgeligne{MPSI/2/1/5}{
  parent = {MPSI/2/1},
  numero = {2.1.5},
  contenu = {
    \cpgeitem{Repérage d’un point dont la trajectoire est connue.}
    \cpgeitem{Vitesse et accélération dans le repère de Frenet pour une trajectoire plane.}},
  capacites = {
    \cpgeitem{Situer qualitativement la direction du vecteur vitesse et du vecteur accélération pour une trajectoire plane.}
    \cpgeitem{Exploiter les liens entre les composantes du vecteur accélération, la courbure de la trajectoire, la norme du vecteur vitesse et sa variation temporelle.}},
  experimentales = {
    \cpgeitem{Réaliser et exploiter quantitativement un enregistrement vidéo d’un mouvement : évolution temporelle des vecteurs vitesse et accélération.}},
  numeriques = {},
}

\cpgerubrique{MPSI/2/2}{
  parent = {MPSI/2},
  numero = {2.2},
  niveau = {2},
  titre = {Lois de Newton},
  chapeau = {
    \cpgeitem{Dans la partie \cpgecettepartie, on cherche d’abord à renforcer les compétences des étudiants relatives à la mise en équations d’un problème, qu’il s’agisse des étapes de bilans de forces ou de projection de la deuxième loi de Newton sur la base choisie. On cherche par ailleurs, sur l’exemple de quelques mouvements simples, à renforcer les compétences d’analyse qualitative d’une équation différentielle : stabilité des solutions, positions d’équilibre, type d’évolution, durée ou période typique d’évolution, etc. Cette pratique s’articule avec l’utilisation d’un langage de programmation pour résoudre des équations différentielles. Enfin, il s’agit aussi de confronter les étudiants aux limites de validité de certains modèles de forces, et ainsi de donner toute leur importance aux étapes de modélisation et de validation d’un modèle.}},
}

\cpgeligne{MPSI/2/2/1}{
  parent = {MPSI/2/2},
  numero = {2.2.1},
  contenu = {
    \cpgeitem{\cpgenotion{Quantité de mouvement}}
    \cpgeitem{Masse d’un système. Conservation de la masse pour système fermé.}},
  capacites = {
    \cpgeitem{Exploiter la conservation de la masse pour un système fermé.}},
  experimentales = {},
  numeriques = {},
}

\cpgeligne{MPSI/2/2/2}{
  parent = {MPSI/2/2},
  numero = {2.2.2},
  contenu = {
    \cpgeitem{Quantité de mouvement d’un point et d’un système de points. Lien avec la vitesse du centre de masse d’un système fermé.}},
  capacites = {
    \cpgeitem{Établir l’expression de la quantité de mouvement pour un système de deux points sous la forme : $\cpgevec p=m\cpgevec v(G)$.}},
  experimentales = {},
  numeriques = {},
}

\cpgeligne{MPSI/2/2/3}{
  parent = {MPSI/2/2},
  numero = {2.2.3},
  contenu = {
    \cpgeitem{Première loi de Newton : principe d’inertie. Référentiels galiléens.}},
  capacites = {
    \cpgeitem{Décrire le mouvement relatif de deux référentiels galiléens.}},
  experimentales = {},
  numeriques = {},
}

\cpgeligne{MPSI/2/2/4}{
  parent = {MPSI/2/2},
  numero = {2.2.4},
  contenu = {
    \cpgeitem{Notion de force. Troisième loi de Newton.}},
  capacites = {
    \cpgeitem{Établir un bilan des forces sur un système ou sur plusieurs systèmes en interaction et en rendre compte sur un schéma.}},
  experimentales = {},
  numeriques = {},
}

\cpgeligne{MPSI/2/2/5}{
  parent = {MPSI/2/2},
  numero = {2.2.5},
  contenu = {
    \cpgeitem{Deuxième loi de Newton.}},
  capacites = {
    \cpgeitem{Déterminer les équations du mouvement d’un point matériel ou du centre de masse d’un système fermé dans un référentiel galiléen.}},
  experimentales = {
    \cpgeitem{Mettre en œuvre un protocole expérimental permettant d’étudier une loi de force par exemple à l’aide d’un microcontrôleur.}},
  numeriques = {},
}

\cpgeligne{MPSI/2/2/6}{
  parent = {MPSI/2/2},
  numero = {2.2.6},
  contenu = {
    \cpgeitem{Force de gravitation.}
    \cpgeitem{Modèle du champ de pesanteur uniforme au voisinage de la surface d’une planète.}
    \cpgeitem{Mouvement dans le champ de pesanteur uniforme.}},
  capacites = {
    \cpgeitem{Etudier le mouvement d’un système modélisé par un point matériel dans un champ de pesanteur uniforme en l’absence de frottement.}},
  experimentales = {},
  numeriques = {},
}

\cpgeligne{MPSI/2/2/7}{
  parent = {MPSI/2/2},
  numero = {2.2.7},
  contenu = {
    \cpgeitem{Modèles d’une force de frottement fluide. Influence de la résistance de l’air sur un mouvement de chute.}},
  capacites = {
    \cpgeitem{Exploiter, sans la résoudre analytiquement, une équation différentielle : analyse en ordres de grandeur, détermination de la vitesse limite, utilisation des résultats obtenus par simulation numérique.}
    \cpgeitem{Écrire une équation adimensionnée.}},
  experimentales = {
    \cpgeitem{Mettre en œuvre un protocole expérimental de mesure de frottements fluides.}},
  numeriques = {},
}

\cpgeligne{MPSI/2/2/8}{
  parent = {MPSI/2/2},
  numero = {2.2.8},
  contenu = {
    \cpgeitem{Tension d’un fil. Pendule simple.}},
  capacites = {
    \cpgeitem{Établir l’équation du mouvement du pendule simple. Justifier l’analogie avec l'oscillateur harmonique dans le cadre de l'approximation linéaire.}},
  experimentales = {},
  numeriques = {},
}

\cpgerubrique{MPSI/2/3}{
  parent = {MPSI/2},
  numero = {2.3},
  niveau = {2},
  titre = {Approche énergétique du mouvement d'un point matériel},
  chapeau = {
    \cpgeitem{La partie \cpgecettepartie vise à construire une démarche alternative et complémentaire pour l’étude d’une situation relevant de la mécanique – et plus généralement de la physique – fondée sur la conservation de certaines grandeurs – ici, l’énergie mécanique. Cette approche est l’occasion d’illustrer la capacité prédictive des analyses graphiques et numériques, par exemple pour pouvoir décrire un comportement à partir d’une représentation graphique de l’énergie potentielle dans le cas d’un mouvement conservatif.}},
}

\cpgeligne{MPSI/2/3/1}{
  parent = {MPSI/2/3},
  numero = {2.3.1},
  contenu = {
    \cpgeitem{\cpgenotion{Puissance, travail et énergie cinétique}}
    \cpgeitem{Puissance  et  travail  d’une  force  dans un référentiel}},
  capacites = {
    \cpgeitem{Reconnaître le caractère moteur ou résistant d'une force.}},
  experimentales = {},
  numeriques = {},
}

\cpgeligne{MPSI/2/3/2}{
  parent = {MPSI/2/3},
  numero = {2.3.2},
  contenu = {
    \cpgeitem{Théorèmes de l’énergie cinétique et de la puissance cinétique dans un référentiel galiléen, dans le cas d’un système modélisé par un point matériel.}},
  capacites = {
    \cpgeitem{Utiliser le théorème approprié en fonction du contexte.}},
  experimentales = {},
  numeriques = {},
}

\cpgeligne{MPSI/2/3/3}{
  parent = {MPSI/2/3},
  numero = {2.3.3},
  contenu = {
    \cpgeitem{\cpgenotion{Champ de force conservative et énergie potentielle}}
    \cpgeitem{Énergie potentielle.}
    \cpgeitem{Lien entre un champ de force conservative et l’énergie potentielle. Gradient.}},
  capacites = {
    \cpgeitem{Établir et citer les expressions de l’énergie potentielle de pesanteur (champ uniforme), de l’énergie potentielle gravitationnelle (champ créé par un astre ponctuel), de l’énergie potentielle élastique.}
    \cpgeitem{Déterminer l’expression d’une force à partir de l’énergie potentielle, l’expression du gradient étant fournie.}
    \cpgeitem{Déduire qualitativement, en un point du graphe d’une fonction énergie potentielle, le sens et l’intensité de la force associée.}},
  experimentales = {},
  numeriques = {},
}

\cpgeligne{MPSI/2/3/4}{
  parent = {MPSI/2/3},
  numero = {2.3.4},
  contenu = {
    \cpgeitem{Énergie mécanique}
    \cpgeitem{Énergie mécanique. Théorème de l’énergie mécanique.}
    \cpgeitem{Mouvement conservatif.}},
  capacites = {
    \cpgeitem{Distinguer force conservative et force non conservative.}
    \cpgeitem{Reconnaître les cas de conservation de l'énergie mécanique.}
    \cpgeitem{Utiliser les conditions initiales.}},
  experimentales = {},
  numeriques = {},
}

\cpgeligne{MPSI/2/3/5}{
  parent = {MPSI/2/3},
  numero = {2.3.5},
  contenu = {
    \cpgeitem{Mouvement conservatif à une dimension.}},
  capacites = {
    \cpgeitem{Identifier sur un graphe d’énergie potentielle une barrière et un puits de potentiel.}
    \cpgeitem{Déduire d’un graphe d’énergie potentielle le comportement qualitatif : trajectoire bornée ou non, mouvement périodique, positions de vitesse nulle.}},
  experimentales = {},
  numeriques = {},
}

\cpgeligne{MPSI/2/3/6}{
  parent = {MPSI/2/3},
  numero = {2.3.6},
  contenu = {
    \cpgeitem{Positions d’équilibre. Stabilité.}},
  capacites = {
    \cpgeitem{Déduire d’un graphe d’énergie potentielle l’existence de positions d’équilibre. Analyser qualitativement la nature, stable ou instable, de ces positions.}},
  experimentales = {},
  numeriques = {},
}

\cpgeligne{MPSI/2/3/7}{
  parent = {MPSI/2/3},
  numero = {2.3.7},
  contenu = {
    \cpgeitem{Petits mouvements au voisinage d’une position d’équilibre stable, approximation locale par un puits de potentiel harmonique.}},
  capacites = {
    \cpgeitem{Établir l’équation différentielle du mouvement au voisinage d’une position d’équilibre.}},
  experimentales = {},
  numeriques = {
    \cpgeitem{à l’aide d’un langage de programmation, résoudre numériquement une équation différentielle du deuxième ordre non-linéaire et faire apparaître l’effet des termes non-linéaires.}},
}

\cpgerubrique{MPSI/2/4}{
  parent = {MPSI/2},
  numero = {2.4},
  niveau = {2},
  titre = {Mouvement de particules chargées dans des champs électrique et magnétostatique, uniformes et stationnaires},
  chapeau = {
    \cpgeitem{La partie \cpgecettepartie introduit l’expression de la force de Lorentz ainsi que deux situations de base sur lesquelles les étudiants doivent être autonomes dans la résolution, attestant en cela de l’acquisition d’une certaine aisance à ce stade de leur formation. Des situations physiques variées sont en capacité d’illustrer concrètement cette partie qui ne doit pas se réduire à des développements calculatoires ou des illustrations graphiques.}},
}

\cpgeligne{MPSI/2/4/1}{
  parent = {MPSI/2/4},
  numero = {2.4.1},
  contenu = {
    \cpgeitem{Force de Lorentz exercée sur une charge ponctuelle ; champs électrique et magnétique.}},
  capacites = {
    \cpgeitem{Évaluer les ordres de grandeur des forces électrique ou magnétique et les comparer à ceux des forces gravitationnelles.}},
  experimentales = {},
  numeriques = {},
}

\cpgeligne{MPSI/2/4/2}{
  parent = {MPSI/2/4},
  numero = {2.4.2},
  contenu = {
    \cpgeitem{Puissance de la force de Lorentz.}},
  capacites = {
    \cpgeitem{Justifier qu’un champ électrique peut modifier l’énergie cinétique d’une particule alors qu’un champ magnétique peut courber la trajectoire sans fournir d’énergie à la particule.}},
  experimentales = {},
  numeriques = {},
}

\cpgeligne{MPSI/2/4/3}{
  parent = {MPSI/2/4},
  numero = {2.4.3},
  contenu = {
    \cpgeitem{Mouvement d’une particule chargée dans un champ électrostatique uniforme.}},
  capacites = {
    \cpgeitem{Mettre en équation le mouvement et le caractériser comme un mouvement à vecteur accélération constant.}
    \cpgeitem{Effectuer un bilan énergétique pour déterminer la valeur de la vitesse d'une particule chargée accélérée par une différence de potentiel.}},
  experimentales = {},
  numeriques = {},
}

\cpgeligne{MPSI/2/4/4}{
  parent = {MPSI/2/4},
  numero = {2.4.4},
  contenu = {
    \cpgeitem{Mouvement d’une particule chargée dans un champ magnétostatique uniforme dans le cas où le vecteur vitesse initial est perpendiculaire au champ magnétostatique.}},
  capacites = {
    \cpgeitem{Déterminer le rayon de la trajectoire et le sens de parcours.}},
  experimentales = {},
  numeriques = {},
}

\cpgerubrique{MPSI/2/5}{
  parent = {MPSI/2},
  numero = {2.5},
  niveau = {2},
  titre = {Moment cinétique},
  chapeau = {
    \cpgeitem{La partie \cpgecettepartie est l’occasion d’introduire les notions de moment cinétique et de moment d’une force. L’un des objectifs visés est que les étudiants disposent de représentations concrètes qui permettent de donner du sens aux grandeurs vectorielles et scalaires utilisées ; c’est notamment pour cela que le bras de levier est introduit. L’accent est mis sur l’identification des situations où le moment cinétique est conservé.}},
}

\cpgeligne{MPSI/2/5/1}{
  parent = {MPSI/2/5},
  numero = {2.5.1},
  contenu = {
    \cpgeitem{Moment cinétique d’un point matériel par rapport à un point et par rapport à un axe orienté.}},
  capacites = {
    \cpgeitem{Relier la direction et le sens du vecteur moment cinétique aux caractéristiques du mouvement.}},
  experimentales = {},
  numeriques = {},
}

\cpgeligne{MPSI/2/5/2}{
  parent = {MPSI/2/5},
  numero = {2.5.2},
  contenu = {
    \cpgeitem{Moment cinétique d’un système discret de points par rapport à un axe orienté.}},
  capacites = {
    \cpgeitem{Utiliser le caractère algébrique du moment cinétique scalaire.}},
  experimentales = {},
  numeriques = {},
}

\cpgeligne{MPSI/2/5/3}{
  parent = {MPSI/2/5},
  numero = {2.5.3},
  contenu = {
    \cpgeitem{Moment d’une force par rapport à un point ou un axe orienté.}},
  capacites = {
    \cpgeitem{Calculer le moment d’une force par rapport à un axe orienté en utilisant le bras de levier.}},
  experimentales = {},
  numeriques = {},
}

\cpgeligne{MPSI/2/5/4}{
  parent = {MPSI/2/5},
  numero = {2.5.4},
  contenu = {
    \cpgeitem{Théorème du moment cinétique en un point fixe dans un référentiel galiléen.}
    \cpgeitem{Conservation du moment cinétique.}},
  capacites = {
    \cpgeitem{Identifier les cas de conservation du moment cinétique.}},
  experimentales = {},
  numeriques = {},
}

\cpgerubrique{MPSI/2/6}{
  parent = {MPSI/2},
  numero = {2.6},
  niveau = {2},
  titre = {Mouvements dans un champ de force centrale conservatif},
  chapeau = {
    \cpgeitem{La partie \cpgecettepartie est notamment motivée par ses nombreuses applications possibles. On discute la nature de la trajectoire sur un graphe donnant l’énergie potentielle effective et, dans le cas d’un champ newtonien (lois de Kepler), on ne poursuit l’étude que dans le cas d'une trajectoire circulaire. Le caractère elliptique des trajectoires associées à un état lié est affirmé sans qu’aucune étude géométrique des ellipses ne soit prévue ; on utilise dans ce cas les constantes du mouvement (moment cinétique et énergie mécanique) pour exprimer l’énergie de la trajectoire elliptique en fonction du demi-grand axe.}},
}

\cpgeligne{MPSI/2/6/1}{
  parent = {MPSI/2/6},
  numero = {2.6.1},
  contenu = {
    \cpgeitem{Point matériel soumis à un champ de force centrale.}},
  capacites = {
    \cpgeitem{Établir la conservation du moment cinétique à partir du théorème du moment cinétique.}
    \cpgeitem{Établir les conséquences de la conservation du moment cinétique : mouvement plan, loi des aires.}},
  experimentales = {},
  numeriques = {},
}

\cpgeligne{MPSI/2/6/2}{
  parent = {MPSI/2/6},
  numero = {2.6.2},
  contenu = {
    \cpgeitem{\cpgenotion{Point matériel soumis à un champ de force centrale conservatif}}
    \cpgeitem{Conservation de l’énergie mécanique.}
    \cpgeitem{Énergie potentielle effective. État lié et état de diffusion.}},
  capacites = {
    \cpgeitem{Exprimer l’énergie mécanique d’un système conservatif ponctuel à partir de l’équation du mouvement.}
    \cpgeitem{Exprimer la conservation de l’énergie mécanique et construire une énergie potentielle effective.}
    \cpgeitem{Décrire qualitativement le mouvement radial à l’aide de l’énergie potentielle effective.}
    \cpgeitem{Relier le caractère borné du mouvement radial à la valeur de l’énergie mécanique.}},
  experimentales = {},
  numeriques = {
    \cpgeitem{à l’aide d’un langage de programmation, obtenir des trajectoires d’un point matériel soumis à un champ de force centrale conservatif.}},
}

\cpgeligne{MPSI/2/6/3}{
  parent = {MPSI/2/6},
  numero = {2.6.3},
  contenu = {
    \cpgeitem{Cas particulier du champ newtonien}
    \cpgeitem{Lois de Kepler.}},
  capacites = {
    \cpgeitem{Énoncer les lois de Kepler pour les planètes et les transposer au cas des satellites terrestres.}},
  experimentales = {},
  numeriques = {},
}

\cpgeligne{MPSI/2/6/4}{
  parent = {MPSI/2/6},
  numero = {2.6.4},
  contenu = {
    \cpgeitem{Cas particulier du mouvement circulaire : satellite, planète.}},
  capacites = {
    \cpgeitem{Établir que le mouvement est uniforme et déterminer sa période.}
    \cpgeitem{Établir la troisième loi de Kepler dans le cas particulier de la trajectoire circulaire. Exploiter sans démonstration sa généralisation au cas d’une trajectoire elliptique.}},
  experimentales = {},
  numeriques = {},
}

\cpgeligne{MPSI/2/6/5}{
  parent = {MPSI/2/6},
  numero = {2.6.5},
  contenu = {
    \cpgeitem{Énergie mécanique dans le cas du mouvement circulaire et dans le cas du mouvement elliptique.}},
  capacites = {
    \cpgeitem{Exprimer l’énergie mécanique pour le mouvement circulaire.}
    \cpgeitem{Exprimer l’énergie mécanique pour le mouvement elliptique en fonction du demi-grand axe.}},
  experimentales = {},
  numeriques = {},
}

\cpgeligne{MPSI/2/6/6}{
  parent = {MPSI/2/6},
  numero = {2.6.6},
  contenu = {
    \cpgeitem{\cpgenotion{Satellites terrestres}}
    \cpgeitem{Satellites géostationnaire, de localisation et de navigation, météorologique.}},
  capacites = {
    \cpgeitem{Différencier les orbites des satellites terrestres en fonction de leurs missions.}
    \cpgeitem{Déterminer l’altitude d’un satellite géostationnaire et justifier sa localisation dans le plan équatorial.}},
  experimentales = {},
  numeriques = {},
}

\cpgerubrique{MPSI/2/7}{
  parent = {MPSI/2},
  numero = {2.7},
  niveau = {2},
  titre = {Mouvement d’un solide},
  chapeau = {
    \cpgeitem{Concernant le solide en rotation autour d’un axe fixe dans la partie \cpgecettepartie, il s’agit de définir le mouvement en remarquant que tout point du solide décrit un cercle autour de l’axe avec une même vitesse angulaire et de déterminer la vitesse de chaque point en fonction de celle-ci et de la distance à l’axe de rotation.}
    \cpgeitem{Des exemples de dynamique du solide sont introduits (translation et rotation autour d’un axe fixe dans un référentiel galiléen), avec toutefois des limitations strictes : l’étude générale d’un mouvement composé d’une translation dans un référentiel galiléen et d’une rotation autour d’un axe fixe dans le référentiel barycentrique ne figure pas au programme. L’étude du mouvement d’un solide en rotation autour d’un axe gardant une direction fixe dans un référentiel galiléen mais pour lequel l’axe de rotation est en mouvement est exclue.}},
}

\cpgeligne{MPSI/2/7/1}{
  parent = {MPSI/2/7},
  numero = {2.7.1},
  contenu = {
    \cpgeitem{\cpgenotion{Description du mouvement d'un solide dans deux cas particuliers}}
    \cpgeitem{Définition d’un solide.}},
  capacites = {
    \cpgeitem{Différencier un solide d’un système déformable.}},
  experimentales = {},
  numeriques = {},
}

\cpgeligne{MPSI/2/7/2}{
  parent = {MPSI/2/7},
  numero = {2.7.2},
  contenu = {
    \cpgeitem{Translation.}},
  capacites = {
    \cpgeitem{Reconnaître et décrire une translation rectiligne ainsi qu’une translation circulaire.}},
  experimentales = {},
  numeriques = {},
}

\cpgeligne{MPSI/2/7/3}{
  parent = {MPSI/2/7},
  numero = {2.7.3},
  contenu = {
    \cpgeitem{Rotation autour d'un axe fixe.}},
  capacites = {
    \cpgeitem{Décrire la trajectoire d’un point quelconque du solide et exprimer sa vitesse en fonction de sa distance à l’axe et de la vitesse angulaire.}},
  experimentales = {},
  numeriques = {},
}

\cpgeligne{MPSI/2/7/4}{
  parent = {MPSI/2/7},
  numero = {2.7.4},
  contenu = {
    \cpgeitem{\cpgenotion{Théorème scalaire du moment cinétique appliqué au solide mobile autour d’un axe fixe}}
    \cpgeitem{Moment cinétique d’un solide en rotation autour d’un axe : moment d’inertie.}},
  capacites = {
    \cpgeitem{Exploiter, pour un solide, la relation entre le moment cinétique scalaire, la vitesse angulaire de rotation et le moment d’inertie fourni.}
    \cpgeitem{Relier qualitativement le moment d’inertie à la répartition des masses.}},
  experimentales = {},
  numeriques = {},
}

\cpgeligne{MPSI/2/7/5}{
  parent = {MPSI/2/7},
  numero = {2.7.5},
  contenu = {
    \cpgeitem{Couple.}},
  capacites = {
    \cpgeitem{Définir un couple.}},
  experimentales = {},
  numeriques = {},
}

\cpgeligne{MPSI/2/7/6}{
  parent = {MPSI/2/7},
  numero = {2.7.6},
  contenu = {
    \cpgeitem{Liaison pivot.}},
  capacites = {
    \cpgeitem{Définir une liaison pivot et justifier le moment qu’elle peut produire.}},
  experimentales = {},
  numeriques = {},
}

\cpgeligne{MPSI/2/7/7}{
  parent = {MPSI/2/7},
  numero = {2.7.7},
  contenu = {
    \cpgeitem{Théorème scalaire du moment cinétique appliqué au solide en rotation autour d’un axe fixe dans un référentiel galiléen.}},
  capacites = {
    \cpgeitem{Exploiter le théorème scalaire du moment cinétique appliqué au solide en rotation autour d’un axe fixe dans un référentiel galiléen.}},
  experimentales = {},
  numeriques = {},
}

\cpgeligne{MPSI/2/7/8}{
  parent = {MPSI/2/7},
  numero = {2.7.8},
  contenu = {
    \cpgeitem{Pendule pesant.}},
  capacites = {
    \cpgeitem{Établir l’équation du mouvement.}
    \cpgeitem{Établir une intégrale première du mouvement.}},
  experimentales = {
    \cpgeitem{Réaliser l’étude énergétique d’un pendule pesant et mettre en évidence une diminution de l’énergie mécanique.}},
  numeriques = {
    \cpgeitem{à l’aide d’un langage de programmation, mettre en évidence le non isochronisme des oscillations.}},
}

\cpgeligne{MPSI/2/7/9}{
  parent = {MPSI/2/7},
  numero = {2.7.9},
  contenu = {
    \cpgeitem{\cpgenotion{Approche énergétique du mouvement d’un solide en rotation autour d’un axe fixe orienté, dans un référentiel galiléen}}
    \cpgeitem{Énergie cinétique d’un solide en rotation autour d’un axe fixe.}},
  capacites = {
    \cpgeitem{Utiliser l’expression de l’énergie cinétique, l’expression du moment d’inertie étant fournie.}},
  experimentales = {},
  numeriques = {},
}

\cpgeligne{MPSI/2/7/10}{
  parent = {MPSI/2/7},
  numero = {2.7.10},
  contenu = {
    \cpgeitem{Théorème de l’énergie cinétique pour un solide en rotation autour d’un axe fixe.}},
  capacites = {
    \cpgeitem{Établir, dans ce cas, l’équivalence entre le théorème scalaire du moment cinétique et celui de l’énergie cinétique.}},
  experimentales = {},
  numeriques = {},
}

\cpgerubrique{MPSI/3}{
  parent = {MPSI},
  numero = {3},
  niveau = {1},
  titre = {L’énergie : conversions et transferts},
  chapeau = {
    \cpgeitem{Après avoir mis l’accent sur le passage fondamental d’une réalité microscopique à des grandeurs mesurables macroscopiques, la partie propose, en s'appuyant sur des exemples concrets, de poursuivre la description et l’étude de la matière à l’échelle macroscopique, et d’aborder les deux principes fondamentaux de la thermodynamique. Les capacités identifiées doivent être introduites en s’appuyant dès que possible sur des dispositifs expérimentaux qui permettent ainsi leur acquisition progressive et authentique.}
    \cpgeitem{On utilise les notations suivantes : pour une grandeur extensive «~$A$~», «~$a$~» sera la grandeur massique associée et «~$A_m$~» la grandeur molaire associée.}},
}

\cpgerubrique{MPSI/3/1}{
  parent = {MPSI/3},
  numero = {3.1},
  niveau = {2},
  titre = {Descriptions microscopique et macroscopique d’un système à l'équilibre},
  chapeau = {},
}

\cpgeligne{MPSI/3/1/1}{
  parent = {MPSI/3/1},
  numero = {3.1.1},
  contenu = {
    \cpgeitem{Échelles microscopique, mésoscopique, et macroscopique. Libre parcours moyen.}},
  capacites = {
    \cpgeitem{Définir l’échelle mésoscopique et en expliquer la nécessité.}
    \cpgeitem{Citer quelques ordres de grandeur de libres parcours moyens.}},
  experimentales = {},
  numeriques = {},
}

\cpgeligne{MPSI/3/1/2}{
  parent = {MPSI/3/1},
  numero = {3.1.2},
  contenu = {
    \cpgeitem{État microscopique et état macroscopique.}},
  capacites = {
    \cpgeitem{Préciser les paramètres nécessaires à la description d’un état microscopique et d’un état macroscopique sur un exemple.}},
  experimentales = {},
  numeriques = {},
}

\cpgeligne{MPSI/3/1/3}{
  parent = {MPSI/3/1},
  numero = {3.1.3},
  contenu = {
    \cpgeitem{Distribution des vitesses moléculaires d'un gaz (homogénéité et isotropie).}
    \cpgeitem{Vitesse quadratique moyenne.}
    \cpgeitem{Température cinétique. Exemple du gaz parfait monoatomique : $E_c=\frac32k_{\mathrm{B}} T$.}},
  capacites = {
    \cpgeitem{Calculer l’ordre de grandeur d’une vitesse quadratique moyenne dans un gaz parfait.}},
  experimentales = {},
  numeriques = {},
}

\cpgeligne{MPSI/3/1/4}{
  parent = {MPSI/3/1},
  numero = {3.1.4},
  contenu = {
    \cpgeitem{Système thermodynamique.}},
  capacites = {
    \cpgeitem{Identifier un système ouvert, un système fermé, un système isolé.}},
  experimentales = {},
  numeriques = {},
}

\cpgeligne{MPSI/3/1/5}{
  parent = {MPSI/3/1},
  numero = {3.1.5},
  contenu = {
    \cpgeitem{État d’équilibre d’un système soumis aux seules forces de pression.}
    \cpgeitem{Pression, température, volume, équation d’état.}
    \cpgeitem{Grandeur extensive, grandeur intensive.}
    \cpgeitem{Exemples du gaz parfait et d’une phase condensée indilatable et incompressible.}},
  capacites = {
    \cpgeitem{Calculer une pression à partir d’une condition d’équilibre mécanique.}
    \cpgeitem{Déduire une température d’une condition d’équilibre thermique.}
    \cpgeitem{Citer quelques ordres de grandeur de volumes molaires ou massiques dans les conditions usuelles de pression et de température.}
    \cpgeitem{Citer et utiliser l’équation d’état des gaz parfaits.}},
  experimentales = {},
  numeriques = {},
}

\cpgeligne{MPSI/3/1/6}{
  parent = {MPSI/3/1},
  numero = {3.1.6},
  contenu = {
    \cpgeitem{Énergie interne d’un système. Capacité thermique à volume constant dans le cas du gaz parfait.}},
  capacites = {
    \cpgeitem{Exprimer l’énergie interne d’un gaz parfait monoatomique à partir de l’interprétation microscopique de la température.}
    \cpgeitem{Exploiter la propriété $U_m=U_m(T)$ pour un gaz parfait.}},
  experimentales = {},
  numeriques = {},
}

\cpgeligne{MPSI/3/1/7}{
  parent = {MPSI/3/1},
  numero = {3.1.7},
  contenu = {
    \cpgeitem{Énergie interne et capacité thermique à volume constant d’une phase condensée considérée incompressible et indilatable.}},
  capacites = {
    \cpgeitem{Exploiter la propriété $U_m=U_m(T)$ pour une phase condensée incompressible et indilatable.}},
  experimentales = {},
  numeriques = {},
}

\cpgeligne{MPSI/3/1/8}{
  parent = {MPSI/3/1},
  numero = {3.1.8},
  contenu = {
    \cpgeitem{Approximation des phases condensées peu compressibles et peu dilatables.}},
  capacites = {
    \cpgeitem{Interpréter graphiquement la différence de compressibilité entre un liquide et un gaz à partir d’isothermes expérimentales.}},
  experimentales = {},
  numeriques = {},
}

\cpgeligne{MPSI/3/1/9}{
  parent = {MPSI/3/1},
  numero = {3.1.9},
  contenu = {
    \cpgeitem{Du gaz réel au gaz parfait.}},
  capacites = {
    \cpgeitem{Comparer le comportement d’un gaz réel au modèle du gaz parfait sur des réseaux d’isothermes expérimentales en coordonnées de Clapeyron ou d’Amagat.}},
  experimentales = {},
  numeriques = {},
}

\cpgeligne{MPSI/3/1/10}{
  parent = {MPSI/3/1},
  numero = {3.1.10},
  contenu = {
    \cpgeitem{Corps pur diphasé en équilibre. Diagramme de phases $(P,T)$.}
    \cpgeitem{Cas de l’équilibre liquide-vapeur : diagramme de Clapeyron $(P,v)$, titre en vapeur.}},
  capacites = {
    \cpgeitem{Analyser un diagramme de phase expérimental $(P,T)$.}
    \cpgeitem{Proposer un jeu de variables d’état suffisant pour caractériser l’état d’équilibre d’un corps pur diphasé soumis aux seules forces de pression.}
    \cpgeitem{Positionner les phases dans les diagrammes $(P,T)$ et $(P,v)$.}
    \cpgeitem{Déterminer la composition d’un mélange diphasé en un point d’un diagramme $(P,v)$.}},
  experimentales = {},
  numeriques = {},
}

\cpgerubrique{MPSI/3/2}{
  parent = {MPSI/3},
  numero = {3.2},
  niveau = {2},
  titre = {Énergie échangée par un système au cours d’une transformation},
  chapeau = {},
}

\cpgeligne{MPSI/3/2/1}{
  parent = {MPSI/3/2},
  numero = {3.2.1},
  contenu = {
    \cpgeitem{Transformation thermodynamique subie par un système.}
    \cpgeitem{Évolutions isochore, isotherme, isobare, monobare, monotherme.}},
  capacites = {
    \cpgeitem{Définir un système adapté à une problématique donnée.}
    \cpgeitem{Exploiter les conditions imposées par le milieu extérieur pour déterminer l’état d’équilibre final.}},
  experimentales = {},
  numeriques = {},
}

\cpgeligne{MPSI/3/2/2}{
  parent = {MPSI/3/2},
  numero = {3.2.2},
  contenu = {
    \cpgeitem{Travail des forces de pression. Transformations isochore, monobare.}},
  capacites = {
    \cpgeitem{Évaluer un travail par découpage en travaux élémentaires et sommation sur un chemin donné dans le cas d’une seule variable.}
    \cpgeitem{Interpréter géométriquement le travail des forces de pression dans un diagramme de Clapeyron.}},
  experimentales = {},
  numeriques = {},
}

\cpgeligne{MPSI/3/2/3}{
  parent = {MPSI/3/2},
  numero = {3.2.3},
  contenu = {
    \cpgeitem{Transferts thermiques.}
    \cpgeitem{Transformation adiabatique.}
    \cpgeitem{Thermostat, transformations monotherme et isotherme.}},
  capacites = {
    \cpgeitem{Distinguer qualitativement les trois types de transferts thermiques : conduction, convection et rayonnement.}
    \cpgeitem{Identifier dans une situation expérimentale le ou les systèmes modélisables par un thermostat.}},
  experimentales = {},
  numeriques = {},
}

\cpgerubrique{MPSI/3/3}{
  parent = {MPSI/3},
  numero = {3.3},
  niveau = {2},
  titre = {Premier principe. Bilans d'énergie},
  chapeau = {
    \cpgeitem{Concernant les bilans d’énergie abordés dans la partie \cpgecettepartie, les expressions des fonctions d’état $U_m(T,V_m)$ et $H_m(T,P)$ sont données si le système ne relève pas du modèle gaz parfait ou du modèle de la phase condensée incompressible et indilatable.}},
}

\cpgeligne{MPSI/3/3/1}{
  parent = {MPSI/3/3},
  numero = {3.3.1},
  contenu = {
    \cpgeitem{Premier principe de la thermodynamique.}},
  capacites = {
    \cpgeitem{Définir un système fermé et établir pour ce système un bilan énergétique faisant intervenir travail et transfert thermique.}
    \cpgeitem{Utiliser le premier principe de la thermodynamique entre deux états voisins.}
    \cpgeitem{Exploiter l’extensivité de l’énergie interne.}
    \cpgeitem{Distinguer le statut de la variation de l’énergie interne du statut des termes d’échange.}
    \cpgeitem{Calculer le transfert thermique sur un chemin donné connaissant le travail et la variation de l’énergie interne.}},
  experimentales = {},
  numeriques = {},
}

\cpgeligne{MPSI/3/3/2}{
  parent = {MPSI/3/3},
  numero = {3.3.2},
  contenu = {
    \cpgeitem{Enthalpie d’un système. Capacité thermique à pression constante dans le cas du gaz parfait et d’une phase condensée incompressible et indilatable.}},
  capacites = {
    \cpgeitem{Exprimer le premier principe sous forme de bilan d’enthalpie dans le cas d’une transformation monobare avec équilibre mécanique dans l’état initial et dans l’état final.}
    \cpgeitem{Exprimer l’enthalpie $H_m(T)$ du gaz parfait à partir de l’énergie interne.}
    \cpgeitem{Justifier que l’enthalpie $H_m$ d’une phase condensée peu compressible et peu dilatable peut être considérée comme une fonction de l’unique variable $T$.}
    \cpgeitem{Citer l’ordre de grandeur de la capacité thermique massique de l’eau liquide.}},
  experimentales = {},
  numeriques = {},
}

\cpgeligne{MPSI/3/3/3}{
  parent = {MPSI/3/3},
  numero = {3.3.3},
  contenu = {
    \cpgeitem{Enthalpie associée à une transition de phase : enthalpie de fusion, enthalpie de vaporisation, enthalpie de sublimation.}},
  capacites = {
    \cpgeitem{Exploiter l’extensivité de l’enthalpie et réaliser des bilans énergétiques en prenant en compte des transitions de phases.}},
  experimentales = {
    \cpgeitem{Mettre en œuvre un protocole expérimental de mesure d'une grandeur thermodynamique énergétique (capacité thermique, enthalpie de fusion, etc.).}},
  numeriques = {},
}

\cpgerubrique{MPSI/3/4}{
  parent = {MPSI/3},
  numero = {3.4},
  niveau = {2},
  titre = {Deuxième principe. Bilans d'entropie},
  chapeau = {
    \cpgeitem{Concernant la partie \cpgecettepartie, l'expression de la fonction d’état entropie est systématiquement donnée et sa construction n’est pas une capacité visée. On cite sans aucun développement quantitatif son interprétation en termes de désordre statistique, de façon à faciliter une interprétation intuitive des bilans d’entropie.}},
}

\cpgeligne{MPSI/3/4/1}{
  parent = {MPSI/3/4},
  numero = {3.4.1},
  contenu = {
    \cpgeitem{Fonction d’état entropie.}},
  capacites = {
    \cpgeitem{Interpréter qualitativement l’entropie en termes de désordre statistique à l’aide de la formule de Boltzmann fournie.}},
  experimentales = {},
  numeriques = {},
}

\cpgeligne{MPSI/3/4/2}{
  parent = {MPSI/3/4},
  numero = {3.4.2},
  contenu = {
    \cpgeitem{Deuxième principe de la thermodynamique : entropie créée, entropie échangée.}
    \cpgeitem{$\Delta S=S_{\text{ech}} + S_{\text{créé}}$ avec $S_{\text{ech}}=\Sigma Q_i/T_i $.}},
  capacites = {
    \cpgeitem{Définir un système fermé et établir pour ce système un bilan entropique.}
    \cpgeitem{Relier la création d’entropie à une ou plusieurs causes physiques de l’irréversibilité.}
    \cpgeitem{Analyser le cas particulier d’un système en évolution adiabatique.}},
  experimentales = {},
  numeriques = {},
}

\cpgeligne{MPSI/3/4/3}{
  parent = {MPSI/3/4},
  numero = {3.4.3},
  contenu = {
    \cpgeitem{Variation d’entropie d’un système.}},
  capacites = {
    \cpgeitem{Utiliser l’expression fournie de la fonction d’état entropie.}
    \cpgeitem{Exploiter l’extensivité de l’entropie.}},
  experimentales = {},
  numeriques = {},
}

\cpgeligne{MPSI/3/4/4}{
  parent = {MPSI/3/4},
  numero = {3.4.4},
  contenu = {
    \cpgeitem{Loi de Laplace.}},
  capacites = {
    \cpgeitem{Citer et utiliser la loi de Laplace et ses conditions d’application.}},
  experimentales = {},
  numeriques = {},
}

\cpgeligne{MPSI/3/4/5}{
  parent = {MPSI/3/4},
  numero = {3.4.5},
  contenu = {
    \cpgeitem{Cas particulier d’une transition de phase.}},
  capacites = {
    \cpgeitem{Citer et utiliser la relation entre les variations d’entropie et d'enthalpie associées à une transition de phase : $\Delta h_{12} (T)=T \Delta s_{12} (T)$}},
  experimentales = {},
  numeriques = {},
}

\cpgerubrique{MPSI/3/5}{
  parent = {MPSI/3},
  numero = {3.5},
  niveau = {2},
  titre = {Machines thermiques},
  chapeau = {
    \cpgeitem{Dans la partie \cpgecettepartie, l'enseignement de la thermodynamique est orienté vers des applications industrielles réelles et motivantes.}},
}

\cpgeligne{MPSI/3/5/1}{
  parent = {MPSI/3/5},
  numero = {3.5.1},
  contenu = {
    \cpgeitem{Application du premier principe et du deuxième principe de la thermodynamique aux machines thermiques cycliques dithermes : rendement, efficacité, théorème de Carnot.}},
  capacites = {
    \cpgeitem{Donner le sens des échanges énergétiques pour un moteur ou un récepteur thermique ditherme.}
    \cpgeitem{Analyser un dispositif concret et le modéliser par une machine cyclique ditherme.}
    \cpgeitem{Définir un rendement ou une efficacité et les relier aux énergies échangées au cours d’un cycle.}
    \cpgeitem{Justifier et utiliser le théorème de Carnot.}
    \cpgeitem{Citer quelques ordres de grandeur des rendements des machines thermiques réelles actuelles.}
    \cpgeitem{Expliquer le principe de la cogénération.}},
  experimentales = {
    \cpgeitem{Mettre en œuvre une machine thermique cyclique ditherme.}},
  numeriques = {},
}

\cpgerubrique{MPSI/4}{
  parent = {MPSI},
  numero = {4},
  niveau = {1},
  titre = {Constitution et transformations de la matière},
  chapeau = {},
}

\cpgerubrique{MPSI/4/1}{
  parent = {MPSI/4},
  numero = {4.1},
  niveau = {2},
  titre = {Transformations de la matière},
  chapeau = {
    \cpgeitem{L’objectif de cette partie est d’amener les étudiants à mobiliser de manière autonome les notions et modèles pour décrire, au niveau macroscopique, un système physico-chimique et son évolution. Il convient que les problématiques abordées, les illustrations et les applications prennent largement appui sur des transformations chimiques rencontrées dans la vie courante, au laboratoire, en milieu industriel ou dans le monde du vivant.}},
}

\cpgerubrique{MPSI/4/1/1}{
  parent = {MPSI/4/1},
  numero = {4.1.1},
  niveau = {3},
  titre = {Description d’un système et de son évolution vers un état final},
  chapeau = {
    \cpgeitem{Les concepts développés dans la partie \cpgecettepartie permettent d’envisager l’optimisation des transformations ou des analyses. L’étude quantitative de l’état final d’un système, siège d’une transformation chimique, est réalisée à partir d’une modélisation par une seule réaction chimique symbolisée par une équation de réaction à laquelle est associée une constante thermodynamique d’équilibre. Il s’agit de prévoir le sens d’évolution de systèmes homogènes ou hétérogènes et de déterminer leur composition dans l’état final.}
    \cpgeitem{Les compétences relatives à cette partie du programme sont ensuite réinvesties au cours de l’année, plus particulièrement au second semestre lors des transformations en solution aqueuse, et en seconde année, notamment dans le cadre de la thermodynamique chimique.}},
}

\cpgeligne{MPSI/4/1/1/1}{
  parent = {MPSI/4/1/1},
  numero = {4.1.1.1},
  contenu = {
    \cpgeitem{\cpgenotion{Système physico-chimique}}
    \cpgeitem{Espèces physico-chimiques.}},
  capacites = {
    \cpgeitem{Recenser les espèces physico-chimiques présentes dans un système.}},
  experimentales = {},
  numeriques = {},
}

\cpgeligne{MPSI/4/1/1/2}{
  parent = {MPSI/4/1/1},
  numero = {4.1.1.2},
  contenu = {
    \cpgeitem{Corps purs et mélanges : concentration en quantité de matière, fraction molaire, pression partielle.}
    \cpgeitem{Composition d’un système physico-chimique}},
  capacites = {
    \cpgeitem{Décrire la composition d’un système à l’aide des grandeurs physiques pertinentes.}},
  experimentales = {},
  numeriques = {},
}

\cpgeligne{MPSI/4/1/1/3}{
  parent = {MPSI/4/1/1},
  numero = {4.1.1.3},
  contenu = {
    \cpgeitem{Variables intensives et extensives.}},
  capacites = {
    \cpgeitem{Identifier le caractère extensif ou intensif d’une variable.}},
  experimentales = {},
  numeriques = {},
}

\cpgeligne{MPSI/4/1/1/4}{
  parent = {MPSI/4/1/1},
  numero = {4.1.1.4},
  contenu = {
    \cpgeitem{\cpgenotion{Transformation chimique d’un système}}
    \cpgeitem{Modélisation d’une transformation par une ou plusieurs réactions chimiques.}},
  capacites = {
    \cpgeitem{Écrire l’équation de la réaction (ou des réactions) qui modélise(nt) une transformation chimique donnée.}},
  experimentales = {},
  numeriques = {},
}

\cpgeligne{MPSI/4/1/1/5}{
  parent = {MPSI/4/1/1},
  numero = {4.1.1.5},
  contenu = {
    \cpgeitem{Équation de réaction ; constante thermodynamique d’équilibre.}},
  capacites = {},
  experimentales = {
    \cpgeitem{Déterminer une constante d’équilibre.}},
  numeriques = {},
}

\cpgeligne{MPSI/4/1/1/6}{
  parent = {MPSI/4/1/1},
  numero = {4.1.1.6},
  contenu = {
    \cpgeitem{Évolution d’un système lors d’une transformation chimique modélisée par une seule réaction chimique : avancement, activité, quotient réactionnel, critère d’évolution.}},
  capacites = {
    \cpgeitem{Décrire qualitativement et quantitativement un système chimique dans l’état initial ou dans un état d’avancement quelconque.}
    \cpgeitem{Exprimer l’activité d’une espèce chimique pure ou dans un mélange dans le cas de solutions aqueuses très diluées ou de mélanges de gaz parfaits avec référence à l’état standard.}
    \cpgeitem{Exprimer le quotient réactionnel.}
    \cpgeitem{Prévoir le sens de l’évolution spontanée d’un système chimique.}},
  experimentales = {},
  numeriques = {},
}

\cpgeligne{MPSI/4/1/1/7}{
  parent = {MPSI/4/1/1},
  numero = {4.1.1.7},
  contenu = {
    \cpgeitem{Composition chimique du système dans l’état final : état d’équilibre chimique, transformation totale.}},
  capacites = {
    \cpgeitem{Identifier un état d’équilibre chimique. Déterminer la composition chimique du système dans l’état final, en distinguant les cas d’équilibre chimique ou de transformation totale, pour une transformation modélisée par une réaction chimique unique.}},
  experimentales = {},
  numeriques = {
    \cpgeitem{déterminer, à l’aide d’un langage de programmation, l’état final d’un système, siège d’une transformation, modélisée par une réaction à partir des conditions initiales et valeur de la constante d’équilibre.}},
}

\cpgerubrique{MPSI/4/1/2}{
  parent = {MPSI/4/1},
  numero = {4.1.2},
  niveau = {3},
  titre = {Évolution temporelle d’un système chimique},
  chapeau = {
    \cpgeitem{La partie \cpgecettepartie permet de dégager expérimentalement les facteurs cinétiques concentration et température. Cette mise en évidence est prolongée par les premières modélisations macroscopiques d’évolution des concentrations avec des lois de vitesse d’ordre simple et d’influence de la température avec la loi d’Arrhenius.}
    \cpgeitem{Les déterminations d’ordre global ou apparent mettent en œuvre la méthode différentielle ou intégrale, et peuvent s’effectuer à l’aide de logiciels dédiés ou de programmes élaborés en langage de programmation, pour l’exploitation des mesures expérimentales dans le cadre d’un réacteur fermé parfaitement agité.}},
}

\cpgeligne{MPSI/4/1/2/1}{
  parent = {MPSI/4/1/2},
  numero = {4.1.2.1},
  contenu = {
    \cpgeitem{\cpgenotion{Cinétique en réacteur fermé de composition uniforme}}
    \cpgeitem{Vitesses de consommation d'un réactif et de formation d'un produit.}
    \cpgeitem{Vitesse de réaction pour une transformation modélisée par une réaction chimique unique supposée sans accumulation d’intermédiaires.}},
  capacites = {
    \cpgeitem{Relier la vitesse de réaction, dans les cas où elle est définie, à la vitesse de consommation d’un réactif ou de formation d’un produit.}},
  experimentales = {},
  numeriques = {},
}

\cpgeligne{MPSI/4/1/2/2}{
  parent = {MPSI/4/1/2},
  numero = {4.1.2.2},
  contenu = {
    \cpgeitem{Lois de vitesse : réactions sans ordre, réactions avec ordre simple (0, 1, 2), ordre global, ordre apparent.}
    \cpgeitem{Temps de demi-vie d’un réactif, temps de demi-réaction.}},
  capacites = {
    \cpgeitem{Exprimer la loi de vitesse si la réaction chimique admet un ordre et déterminer la valeur de la constante cinétique à une température donnée. Déterminer la vitesse de réaction à différentes dates en utilisant une méthode numérique ou graphique. Déterminer un ordre de réaction à l’aide de la méthode différentielle ou à l’aide des temps de demi-réaction.}
    \cpgeitem{Confirmer la valeur d'un ordre par la méthode intégrale, en se limitant strictement à une décomposition d'ordre 0, 1 ou 2 d'un unique réactif, ou se ramenant à un tel cas par dégénérescence de l'ordre ou conditions initiales stœchiométriques.}},
  experimentales = {
    \cpgeitem{Établir une loi de vitesse à partir du suivi temporel d’une grandeur physique.}},
  numeriques = {},
}

\cpgeligne{MPSI/4/1/2/3}{
  parent = {MPSI/4/1/2},
  numero = {4.1.2.3},
  contenu = {
    \cpgeitem{Loi d'Arrhenius ; énergie d'activation.}},
  capacites = {
    \cpgeitem{Déterminer la valeur de l’énergie d’activation d’une réaction chimique à partir de valeurs de la constante cinétique à différentes températures.}},
  experimentales = {
    \cpgeitem{Déterminer l’énergie d’activation d’une réaction chimique.}},
  numeriques = {},
}

\cpgerubrique{MPSI/4/2}{
  parent = {MPSI/4},
  numero = {4.2},
  niveau = {2},
  titre = {Relations entre la structure des entités chimiques et les propriétés physiques macroscopiques},
  chapeau = {
    \cpgeitem{Décrivant la matière au niveau macroscopique par des espèces chimiques aux propriétés physiques et chimiques caractéristiques, le chimiste la modélise au niveau microscopique par des entités chimiques dont les structures électroniques et géométriques permettent d’interpréter et de prévoir ces propriétés.}},
}

\cpgerubrique{MPSI/4/2/1}{
  parent = {MPSI/4/2},
  numero = {4.2.1},
  niveau = {3},
  titre = {Structure des entités chimiques},
  chapeau = {
    \cpgeitem{La partie \cpgecettepartie aborde l’étude de la constitution de la matière au niveau microscopique en s’appuyant sur le tableau périodique des éléments, outil essentiel du chimiste, dans l’objectif de développer progressivement les compétences relatives à l’utilisation des informations qu’il contient pour prévoir, dans cette partie, le nombre de liaisons d’un atome et la nature (polaire, ionique) des liaisons chimiques.}},
}

\cpgeligne{MPSI/4/2/1/1}{
  parent = {MPSI/4/2/1},
  numero = {4.2.1.1},
  contenu = {
    \cpgeitem{\cpgenotion{Modèle de la liaison covalente}}
    \cpgeitem{Liaison covalente localisée.}
    \cpgeitem{Schéma de Lewis d’une molécule ou d’un ion monoatomique ou d’un ion polyatomique pour les éléments des blocs s et p.}},
  capacites = {
    \cpgeitem{Citer les ordres de grandeur de longueurs et d’énergies de liaisons covalentes.}
    \cpgeitem{Déterminer, pour les éléments des blocs s et p, le nombre d’électrons de valence d’un atome à partir de la position de l’élément dans le tableau périodique.}
    \cpgeitem{Établir un schéma de Lewis pertinent pour une molécule ou un ion.}
    \cpgeitem{Identifier les écarts à la règle de l’octet.}},
  experimentales = {},
  numeriques = {},
}

\cpgeligne{MPSI/4/2/1/2}{
  parent = {MPSI/4/2/1},
  numero = {4.2.1.2},
  contenu = {
    \cpgeitem{\cpgenotion{Géométrie et polarité des entités chimiques}}
    \cpgeitem{Électronégativité : liaison polarisée, moment dipolaire, molécule polaire.}},
  capacites = {
    \cpgeitem{Associer qualitativement la géométrie d’une entité à une minimisation de son énergie.}
    \cpgeitem{Comparer les électronégativités de deux atomes à partir de données ou de leurs positions dans le tableau périodique.}
    \cpgeitem{Prévoir la polarisation d’une liaison à partir des électronégativités comparées des deux atomes mis en jeu.}
    \cpgeitem{Relier l’existence ou non d’un moment dipolaire permanent à la structure géométrique donnée d’une molécule.}
    \cpgeitem{Déterminer direction et sens du vecteur moment dipolaire d’une liaison ou d’une molécule de géométrie donnée.}},
  experimentales = {},
  numeriques = {},
}

\cpgerubrique{MPSI/4/2/2}{
  parent = {MPSI/4/2},
  numero = {4.2.2},
  niveau = {3},
  titre = {Relations structure des entités - propriétés physiques macroscopiques},
  chapeau = {
    \cpgeitem{La partie \cpgecettepartie a pour objectif de permettre l’identification des interactions entre entités moléculaires ou ioniques afin d’interpréter, de prévoir ou de comparer certaines propriétés physiques : température de changement d’état, miscibilité, solubilité.}
    \cpgeitem{De nombreuses illustrations et applications dans la vie courante, au niveau du laboratoire ou dans le domaine du vivant peuvent être proposées.}},
}

\cpgeligne{MPSI/4/2/2/1}{
  parent = {MPSI/4/2/2},
  numero = {4.2.2.1},
  contenu = {
    \cpgeitem{\cpgenotion{Interaction entre entités}}
    \cpgeitem{Interactions de van der Waals.}
    \cpgeitem{Liaison hydrogène ou interaction par pont hydrogène.}},
  capacites = {
    \cpgeitem{Citer les ordres de grandeur énergétiques des interactions de van der Waals et de liaisons hydrogène.}
    \cpgeitem{Interpréter l’évolution de températures de changement d’état de corps purs moléculaires à l’aide de l’existence d’interactions de van der Waals ou par pont hydrogène.}},
  experimentales = {},
  numeriques = {},
}

\cpgeligne{MPSI/4/2/2/2}{
  parent = {MPSI/4/2/2},
  numero = {4.2.2.2},
  contenu = {
    \cpgeitem{\cpgenotion{Solubilité ; miscibilité.}}
    \cpgeitem{Grandeurs caractéristiques et propriétés de solvants moléculaires : moment dipolaire, permittivité relative, caractère protogène.}
    \cpgeitem{Mise en solution d’une espèce chimique moléculaire ou ionique.}},
  capacites = {
    \cpgeitem{Associer une propriété d’un solvant moléculaire à une ou des grandeurs caractéristiques.}
    \cpgeitem{Interpréter la miscibilité ou la non-miscibilité de deux solvants.}
    \cpgeitem{Interpréter la solubilité d’une espèce chimique moléculaire ou ionique.}},
  experimentales = {},
  numeriques = {},
}

\cpgerubrique{MPSI/4/3}{
  parent = {MPSI/4},
  numero = {4.3},
  niveau = {2},
  titre = {Structure et propriétés physiques des solides},
  chapeau = {
    \cpgeitem{Les modèles de description microscopique des solides sont présentés dans la partie \cpgecettepartie à partir de l’observation et des propriétés macroscopiques de différents solides cristallisés que l’enseignant est libre de choisir. L’introduction du modèle du cristal parfait se fait sur l’exemple de la maille cubique à faces centrées (CFC), seule maille dont la connaissance est exigible ; l’ensemble des notions associées à cette première étude est réinvesti pour étudier d’autres structures cristallines dont la constitution est alors fournie.}
    \cpgeitem{L’objectif principal de l’étude des cristaux métalliques, covalents et ioniques est d’aborder une nouvelle fois la notion de modèle : les allers-retours entre le niveau macroscopique (solides de différentes natures) et la modélisation microscopique (cristal parfait) permettent de montrer les limites du modèle du cristal parfait et de confronter les prédictions faites par ce modèle aux valeurs expérimentales mesurées sur le solide réel (distances internucléaires et interatomiques, masse volumique, etc.). Ce chapitre constitue une occasion de revenir sur les positions relatives des éléments dans le tableau périodique, en lien avec la nature des interactions assurant la cohésion des édifices présentés, ainsi que sur les interactions intermoléculaires et la notion de solubilisation pour les solides ioniques et moléculaires.}
    \cpgeitem{Une réflexion sur les modèles conduisant à la détermination des différents types de rayons à partir des méthodes expérimentales d’analyse des structures des solides peut être proposée.}},
}

\cpgeligne{MPSI/4/3/1}{
  parent = {MPSI/4/3},
  numero = {4.3.1},
  contenu = {
    \cpgeitem{\cpgenotion{Modèle du cristal parfait}}
    \cpgeitem{Solide amorphe, solide cristallin, solide semi-cristallin ; variétés allotropiques.}},
  capacites = {},
  experimentales = {
    \cpgeitem{Illustrer l’influence des conditions expérimentales sur la formation de solides et de solides cristallins.}},
  numeriques = {},
}

\cpgeligne{MPSI/4/3/2}{
  parent = {MPSI/4/3},
  numero = {4.3.2},
  contenu = {
    \cpgeitem{Description du cristal parfait ; population, coordinence, compacité, masse volumique.}
    \cpgeitem{Rayons métallique, covalent, de van der Waals ou ionique.}},
  capacites = {
    \cpgeitem{Décrire un cristal parfait comme un assemblage de mailles parallélépipédiques.}
    \cpgeitem{Déterminer la population, la coordinence et la compacité pour une structure fournie.}
    \cpgeitem{Déterminer la valeur de la masse volumique d’un matériau cristallisé selon une structure cristalline fournie.}
    \cpgeitem{Relier le rayon métallique, covalent, de van der Waals ou ionique, selon le cas, aux paramètres d’une maille donnée.}},
  experimentales = {
    \cpgeitem{Utiliser un logiciel ou des modèles cristallins pour visualiser des mailles et des sites interstitiels et pour déterminer des paramètres géométriques.}},
  numeriques = {},
}

\cpgeligne{MPSI/4/3/3}{
  parent = {MPSI/4/3},
  numero = {4.3.3},
  contenu = {
    \cpgeitem{Description des modèles d’empilement compact de sphères identiques.}},
  capacites = {
    \cpgeitem{Localiser les interstices tétraédriques et octaédriques entre les plans d’empilement.}},
  experimentales = {},
  numeriques = {},
}

\cpgeligne{MPSI/4/3/4}{
  parent = {MPSI/4/3},
  numero = {4.3.4},
  contenu = {
    \cpgeitem{Maille conventionnelle CFC et ses sites interstitiels.}},
  capacites = {
    \cpgeitem{Localiser, dénombrer les sites tétraédriques et octaédriques d’une maille CFC et déterminer leur habitabilité.}},
  experimentales = {},
  numeriques = {},
}

\cpgeligne{MPSI/4/3/5}{
  parent = {MPSI/4/3},
  numero = {4.3.5},
  contenu = {
    \cpgeitem{Limites du modèle du cristal parfait.}},
  capacites = {
    \cpgeitem{Confronter des données expérimentales aux prévisions du modèle.}},
  experimentales = {},
  numeriques = {},
}

\cpgeligne{MPSI/4/3/6}{
  parent = {MPSI/4/3},
  numero = {4.3.6},
  contenu = {
    \cpgeitem{\cpgenotion{Métaux}}
    \cpgeitem{Cohésion et propriétés physiques des métaux.}},
  capacites = {
    \cpgeitem{Positionner dans le tableau périodique et reconnaître les métaux et non métaux.}
    \cpgeitem{Relier les caractéristiques de la liaison métallique (ordre de grandeur énergétique, non directionnalité) aux propriétés macroscopiques des métaux.}},
  experimentales = {},
  numeriques = {},
}

\cpgeligne{MPSI/4/3/7}{
  parent = {MPSI/4/3},
  numero = {4.3.7},
  contenu = {
    \cpgeitem{\cpgenotion{Solides covalents et moléculaires}}
    \cpgeitem{Cohésion et propriétés physiques des solides covalents et moléculaires.}},
  capacites = {
    \cpgeitem{Relier les caractéristiques des liaisons covalentes, des interactions de van der Waals et des interactions par pont hydrogène (directionnalité ou non, ordre de grandeur des énergies mises en jeu) et les propriétés macroscopiques des solides correspondants.}},
  experimentales = {},
  numeriques = {},
}

\cpgeligne{MPSI/4/3/8}{
  parent = {MPSI/4/3},
  numero = {4.3.8},
  contenu = {
    \cpgeitem{\cpgenotion{Solides ioniques}}
    \cpgeitem{Cohésion et propriétés physiques des solides ioniques.}},
  capacites = {
    \cpgeitem{Relier les caractéristiques de l’interaction ionique dans le cadre du modèle du solide ionique parfait (ordre de grandeur de l’énergie d’interaction, non directionnalité, charge localisée) avec les propriétés macroscopiques des solides ioniques.}},
  experimentales = {},
  numeriques = {},
}

\cpgerubrique{MPSI/4/4}{
  parent = {MPSI/4},
  numero = {4.4},
  niveau = {2},
  titre = {Transformations chimiques en solution aqueuse},
  chapeau = {
    \cpgeitem{Les transformations chimiques en solution aqueuse jouent un rôle essentiel en chimie, en biochimie, dans le domaine du vivant et dans les procédés industriels. Un nombre considérable de développements technologiques et d’analyses environnementales (traitement des eaux, méthodes d’analyse, extraction d’ions métalliques des minerais, générateurs électrochimiques, lutte contre la corrosion, etc) repose sur des transformations acido-basiques, de solubilisation-précipitation et d’oxydo-réduction en solution aqueuse dont la maîtrise est importante pour prévoir, interpréter et optimiser les phénomènes mis en jeu.}
    \cpgeitem{L’objectif de cette partie est donc de présenter différents types de réactions susceptibles d’intervenir en solution aqueuse, d’en déduire des diagrammes de prédominance ou d’existence d’espèces chimiques, notamment des diagrammes potentiel-pH, et de les utiliser comme outil de prévision et d’interprétation des transformations chimiques quel que soit le milieu donné.}
    \cpgeitem{Les conventions de tracé seront toujours précisées. Les choix pédagogiques relatifs au contenu des séances de travail expérimental permettront de contextualiser ces enseignements. Les dosages par titrage sont étudiés exclusivement en travaux pratiques. L’analyse des conditions choisies ou la réflexion conduisant à une proposition de protocole expérimental pour atteindre un objectif donné constituent des mises en situation des enseignements évoqués précédemment. Ces séances de travail expérimental constituent une nouvelle occasion d’aborder qualité et précision de la mesure.}},
}

\cpgerubrique{MPSI/4/4/1}{
  parent = {MPSI/4/4},
  numero = {4.4.1},
  niveau = {3},
  titre = {Réactions acide-base et de précipitation},
  chapeau = {
    \cpgeitem{Les différentes transformations en solution aqueuse abordées dans la partie \cpgecettepartie constituent des illustrations de l’évolution des systèmes chimiques introduites au premier semestre, les étudiants étant amenés à déterminer l’état final d’un système en transformation chimique modélisée par une seule réaction chimique. On montrera qu’il est ainsi possible d’analyser et de simplifier une situation complexe pour parvenir à la décrire rigoureusement et quantitativement, en l’occurrence dans le cas des solutions aqueuses par une seule réaction. Il est cependant important de noter qu’on évite tout calcul inutile de concentration, en privilégiant l’utilisation des diagrammes pour valider le choix de la réaction mise en jeu. Dans ce cadre, aucune formule de calcul de pH n’est exigible.}},
}

\cpgeligne{MPSI/4/4/1/1}{
  parent = {MPSI/4/4/1},
  numero = {4.4.1.1},
  contenu = {
    \cpgeitem{\cpgenotion{Réactions acido-basiques} \begin{itemize} \item constante d’acidité ; \item diagramme de prédominance, de distribution ; \item exemples usuels d’acides et bases : nom, formule et nature – faible ou forte – des acides sulfurique, nitrique, chlorhydrique, phosphorique, acétique, de la soude, l’ion hydrogénocarbonate, l’ammoniac. \end{itemize}}
    \cpgeitem{\cpgenotion{Réactions de dissolution ou de précipitation} \begin{itemize} \item constante de l’équation de dissolution, produit de solubilité $K_s$ ; \item solubilité et condition de précipitation ; \item domaine d’existence ; \item facteurs influençant la solubilité. \end{itemize}}},
  capacites = {
    \cpgeitem{Identifier le caractère acido-basique d’une réaction en solution aqueuse.}
    \cpgeitem{Écrire l’équation de la réaction modélisant une transformation en solution aqueuse en tenant compte des caractéristiques du milieu réactionnel}
    \cpgeitem{(nature des espèces chimiques en présence, pH…) et des observations expérimentales.}
    \cpgeitem{Déterminer la valeur de la constante d’équilibre pour une équation de réaction, combinaison linéaire d’équations dont les constantes thermodynamiques sont connues.}
    \cpgeitem{Déterminer la composition chimique du système dans l’état final, en distinguant les cas d’équilibre chimique et de transformation totale, pour une transformation modélisée par une réaction chimique unique.}
    \cpgeitem{Prévoir l’état de saturation ou de non saturation d’une solution.}
    \cpgeitem{Utiliser les diagrammes de prédominance ou d’existence pour prévoir les espèces incompatibles ou la nature des espèces majoritaires.}
    \cpgeitem{Exploiter des courbes d’évolution de la solubilité d’un solide en fonction d’une variable.}},
  experimentales = {
    \cpgeitem{Mettre en œuvre une réaction acide-base et une réaction de précipitation pour réaliser une analyse quantitative en solution aqueuse.}
    \cpgeitem{Illustrer un procédé de retraitement, de recyclage, de séparation en solution aqueuse.}},
  numeriques = {},
}

\cpgerubrique{MPSI/4/4/2}{
  parent = {MPSI/4/4},
  numero = {4.4.2},
  niveau = {3},
  titre = {Réactions d’oxydo-réduction},
  chapeau = {
    \cpgeitem{L’analyse de transformations mettant en jeu des oxydants et réducteurs usuels et des piles permettent d’aborder, dans la partie \cpgecettepartie les différents concepts associés aux phénomènes d’oxydo-réduction en solution aqueuse. La relation de Nernst (admise en première année) ainsi que la relation entre la constante thermodynamique d'équilibre d’une réaction d’oxydo-réduction et les potentiels standard permettent de prévoir l’évolution des systèmes et le caractère favorisé des transformations.}
    \cpgeitem{Afin de pouvoir étudier l’influence du milieu sur les espèces oxydantes ou réductrices présentes, les acquis sur les réactions acido-basiques et de précipitation-solubilisation en solution aqueuse sont réinvestis.}
    \cpgeitem{Enfin, les diagrammes potentiel-pH sont présentés puis superposés pour prévoir ou interpréter thermodynamiquement des transformations chimiques ; la confrontation avec la réalité amenant à aborder éventuellement des blocages cinétiques en lien avec l’évolution temporelle des systèmes étudiée au premier semestre.}},
}

\cpgeligne{MPSI/4/4/2/1}{
  parent = {MPSI/4/4/2},
  numero = {4.4.2.1},
  contenu = {
    \cpgeitem{\cpgenotion{Oxydants et réducteurs, réactions d’oxydo-réduction}}
    \cpgeitem{Nombre d’oxydation.}
    \cpgeitem{Exemples d’oxydants et de réducteurs minéraux usuels : nom, nature et formule des ions thiosulfate, permanganate, hypochlorite, du peroxyde d’hydrogène.}},
  capacites = {
    \cpgeitem{Relier la position d’un élément dans le tableau périodique et le caractère oxydant ou réducteur du corps simple correspondant.}
    \cpgeitem{Prévoir les nombres d’oxydation extrêmes d’un élément à partir de sa position dans le tableau périodique.}
    \cpgeitem{Identifier l’oxydant et le réducteur d’un couple.}},
  experimentales = {},
  numeriques = {},
}

\cpgeligne{MPSI/4/4/2/2}{
  parent = {MPSI/4/4/2},
  numero = {4.4.2.2},
  contenu = {
    \cpgeitem{Pile, tension à vide, potentiel d’électrode, formule de Nernst, électrodes de référence.}},
  capacites = {
    \cpgeitem{Décrire le fonctionnement d’une pile à partir d’une mesure de tension à vide ou à partir des potentiels d’électrode.}},
  experimentales = {},
  numeriques = {},
}

\cpgeligne{MPSI/4/4/2/3}{
  parent = {MPSI/4/4/2},
  numero = {4.4.2.3},
  contenu = {
    \cpgeitem{Diagrammes de prédominance ou d’existence.}},
  capacites = {
    \cpgeitem{Utiliser les diagrammes de prédominance ou d’existence pour prévoir les espèces incompatibles ou la nature des espèces majoritaires.}},
  experimentales = {},
  numeriques = {},
}

\cpgeligne{MPSI/4/4/2/4}{
  parent = {MPSI/4/4/2},
  numero = {4.4.2.4},
  contenu = {
    \cpgeitem{Aspect thermodynamique des réactions d’oxydo-réduction.}
    \cpgeitem{Dismutation et médiamutation.}},
  capacites = {
    \cpgeitem{Prévoir qualitativement ou quantitativement le caractère thermodynamiquement favorisé ou défavorisé d’une réaction d’oxydo-réduction à partir des potentiels standard des couples.}},
  experimentales = {
    \cpgeitem{Mettre en œuvre une réaction d’oxydo-réduction pour réaliser une analyse quantitative en solution aqueuse.}},
  numeriques = {},
}

\cpgeligne{MPSI/4/4/2/5}{
  parent = {MPSI/4/4/2},
  numero = {4.4.2.5},
  contenu = {
    \cpgeitem{\cpgenotion{Diagrammes potentiel-pH}}
    \cpgeitem{Principe de construction, lecture et utilisation d’un diagramme potentiel-pH.}},
  capacites = {
    \cpgeitem{Identifier les différents domaines d’un diagramme fourni associés à des espèces chimiques données.}
    \cpgeitem{Déterminer la valeur de la pente d’une frontière dans un diagramme potentiel-pH.}
    \cpgeitem{Justifier la position d’une frontière verticale.}
    \cpgeitem{Prévoir le caractère thermodynamiquement favorisé ou non d’une transformation par superposition de diagrammes.}},
  experimentales = {
    \cpgeitem{Réaliser une pile et étudier son fonctionnement.}},
  numeriques = {},
}

\cpgeligne{MPSI/4/4/2/6}{
  parent = {MPSI/4/4/2},
  numero = {4.4.2.6},
  contenu = {
    \cpgeitem{Diagramme potentiel-pH de l’eau}},
  capacites = {
    \cpgeitem{Prévoir la stabilité des espèces dans l’eau. Prévoir une dismutation ou médiamutation en fonction du pH du milieu.}
    \cpgeitem{Confronter les prévisions à des données expérimentales et interpréter d’éventuels écarts en termes cinétiques.}},
  experimentales = {
    \cpgeitem{Mettre en œuvre des réactions d’oxydo-réduction en s’appuyant sur l’utilisation de diagrammes potentiel-pH.}},
  numeriques = {},
}

\cpgerubrique{MPSIA/A1}{
  parent = {MPSI},
  numero = {A1},
  niveau = {1},
  titre = {Matériel},
  chapeau = {
    \cpgeitem{La liste ci-dessous regroupe le matériel que les étudiants doivent savoir utiliser avec l’aide d’une notice simplifiée fournie sous forme de version papier ou numérique. Une utilisation de matériel hors de cette liste lors d’épreuves d’évaluation n’est pas exclue, mais elle doit obligatoirement s’accompagner d’une introduction guidée suffisamment détaillée.}},
}

\cpgeligne{MPSIA/A1/1}{
  parent = {MPSIA/A1},
  numero = {A1.1},
  contenu = {
    \cpgeitem{Domaine optique}},
  capacites = {
    \cpgeitem{\begin{itemize} \item Goniomètre \item Viseur à frontale fixe \item Lunette auto-collimatrice \item Spectromètre à fibre optique \item Laser à gaz \item Lampes spectrales \item Source de lumière blanche à condenseur \end{itemize}}},
  experimentales = {},
  numeriques = {},
}

\cpgeligne{MPSIA/A1/2}{
  parent = {MPSIA/A1},
  numero = {A1.2},
  contenu = {
    \cpgeitem{Domaine électrique}},
  capacites = {
    \cpgeitem{\begin{itemize} \item Oscilloscope numérique \item Carte d’acquisition et logiciel dédié \item Générateur de signaux Basse Fréquence \item Multimètre numérique \item Multiplieur analogique \item Émetteur et récepteur acoustique (domaine audible et domaine ultrasonore) \item Microcontrôleur \end{itemize}}},
  experimentales = {},
  numeriques = {},
}

\cpgeligne{MPSIA/A1/3}{
  parent = {MPSIA/A1},
  numero = {A1.3},
  contenu = {
    \cpgeitem{Domaines mécanique et thermodynamique}},
  capacites = {
    \cpgeitem{\begin{itemize} \item Dynamomètre \item Capteur de pression \item Accéléromètre \item Stroboscope \item Webcam avec logiciel dédié \item Appareil photo numérique ou caméra numérique \item Thermomètre, thermocouple, thermistance, capteur infra-rouge \item Calorimètre \item Machines thermiques dithermes \end{itemize}}},
  experimentales = {},
  numeriques = {},
}

\cpgeligne{MPSIA/A1/4}{
  parent = {MPSIA/A1},
  numero = {A1.4},
  contenu = {
    \cpgeitem{Domaine constitution et transformations de la matière}},
  capacites = {
    \cpgeitem{\begin{itemize} \item Verrerie classique de chimie analytique : burettes, pipettes jaugées et graduées, fioles jaugées, erlenmeyers, béchers, etc. \item Matériel classique du laboratoire de chimie : dispositifs de chauffage ou de refroidissement (bain-marie, bain froid, etc.), dispositifs d’agitation, matériel de filtration sous pression atmosphérique et sous pression réduite \item Spectrophotomètre UV-visible \item pH-mètre et électrodes de mesure \item Voltmètre et électrodes \item Conductimètre et cellule de mesure \item Thermomètre \item Balance de précision \end{itemize}}},
  experimentales = {},
  numeriques = {},
}

\cpgerubrique{MPSIA/A2}{
  parent = {MPSI},
  numero = {A2},
  niveau = {1},
  titre = {Outils mathématiques},
  chapeau = {
    \cpgeitem{L’utilisation d’outils mathématiques est indispensable en physique comme en chimie.}
    \cpgeitem{La capacité à mettre en œuvre de manière autonome certains de ces outils mathématiques dans le cadre des activités relevant de la physique-chimie fait partie des compétences exigibles à la fin de la première année. Le tableau ci-dessous explicite ces outils ainsi que le niveau de maîtrise attendu en fin de première année. Il est complété dans le programme de seconde année.}
    \cpgeitem{Cependant les situations dont la gestion manuelle ne relèverait que de la technicité sont traitées à l’aide d’outils numériques (calculatrices, logiciels de calcul numérique).}},
}

\cpgerubrique{MPSIA/A2/1}{
  parent = {MPSIA/A2},
  numero = {A2.1},
  niveau = {2},
  titre = {Équations algébriques},
  chapeau = {},
}

\cpgeligne{MPSIA/A2/1/1}{
  parent = {MPSIA/A2/1},
  numero = {A2.1.1},
  contenu = {
    \cpgeitem{Systèmes linéaires de $n$ équations à $p$ inconnues.}},
  capacites = {
    \cpgeitem{Identifier les variables (inconnues) nécessaires à la modélisation du problème sous forme d’un système d’équations linéaires.}
    \cpgeitem{Donner l’expression formelle des solutions dans le seul cas $n = p = 2$.}},
  experimentales = {},
  numeriques = {},
}

\cpgeligne{MPSIA/A2/1/2}{
  parent = {MPSIA/A2/1},
  numero = {A2.1.2},
  contenu = {
    \cpgeitem{Équations non linéaires.}},
  capacites = {
    \cpgeitem{Représenter graphiquement une équation de la forme $f(x) = g(x)$.}
    \cpgeitem{Interpréter graphiquement la ou les solutions.}},
  experimentales = {},
  numeriques = {},
}

\cpgerubrique{MPSIA/A2/2}{
  parent = {MPSIA/A2},
  numero = {A2.2},
  niveau = {2},
  titre = {Équations différentielles},
  chapeau = {},
}

\cpgeligne{MPSIA/A2/2/1}{
  parent = {MPSIA/A2/2},
  numero = {A2.2.1},
  contenu = {
    \cpgeitem{Équations différentielles linéaires à coefficients constants.}},
  capacites = {
    \cpgeitem{Identifier l’ordre.}
    \cpgeitem{Mettre l’équation sous forme canonique.}},
  experimentales = {},
  numeriques = {},
}

\cpgeligne{MPSIA/A2/2/2}{
  parent = {MPSIA/A2/2},
  numero = {A2.2.2},
  contenu = {
    \cpgeitem{Équations différentielles linéaires du premier ordre à coefficients constants : $y' + ay = f(x)$}},
  capacites = {
    \cpgeitem{Trouver la solution générale de l’équation sans second membre (équation homogène).}
    \cpgeitem{Trouver l’expression des solutions lorsque $f(x)$ est constante ou de la forme $A.cos(\omega x+\varphi)$ (en utilisant la notation complexe).}},
  experimentales = {},
  numeriques = {},
}

\cpgeligne{MPSIA/A2/2/3}{
  parent = {MPSIA/A2/2},
  numero = {A2.2.3},
  contenu = {
    \cpgeitem{Équations différentielles linéaires du deuxième ordre à coefficients constants : $y'' + ay' + by = f(x)$.}},
  capacites = {
    \cpgeitem{Utiliser l’équation caractéristique pour trouver la solution générale de l’équation sans second membre.}
    \cpgeitem{Prévoir le caractère borné ou non de ses solutions (critère de stabilité).}
    \cpgeitem{Trouver l’expression des solutions lorsque $f(x)$ est constante ou de la forme $A.exp(\lambda x)$ avec $\lambda$ complexe.}
    \cpgeitem{Trouver la solution de l’équation complète correspondant à des conditions initiales données.}
    \cpgeitem{Représenter graphiquement cette solution.}},
  experimentales = {},
  numeriques = {},
}

\cpgeligne{MPSIA/A2/2/4}{
  parent = {MPSIA/A2/2},
  numero = {A2.2.4},
  contenu = {
    \cpgeitem{Autres équations différentielles d’ordre 1 ou 2.}},
  capacites = {
    \cpgeitem{Obtenir une intégrale première d’une équation de Newton $x'' = f(x)$ et l’exploiter graphiquement.}
    \cpgeitem{Séparer les variables d’une équation du premier ordre à variables séparables.}
    \cpgeitem{Faire le lien entre les conditions initiales et le graphe de la solution correspondante.}},
  experimentales = {},
  numeriques = {},
}

\cpgerubrique{MPSIA/A2/3}{
  parent = {MPSIA/A2},
  numero = {A2.3},
  niveau = {2},
  titre = {Fonctions},
  chapeau = {},
}

\cpgeligne{MPSIA/A2/3/1}{
  parent = {MPSIA/A2/3},
  numero = {A2.3.1},
  contenu = {
    \cpgeitem{Fonctions usuelles.}},
  capacites = {
    \cpgeitem{Exponentielle, logarithme népérien et décimal, cosinus, sinus, tangente, puissance réelle $(x\mapsto x^a)$.}},
  experimentales = {},
  numeriques = {},
}

\cpgeligne{MPSIA/A2/3/2}{
  parent = {MPSIA/A2/3},
  numero = {A2.3.2},
  contenu = {
    \cpgeitem{Dérivée.}
    \cpgeitem{Notation $\cpgedd x/\cpgedd t$.}},
  capacites = {
    \cpgeitem{Utiliser la formule de Taylor à l’ordre un ou deux ; interpréter graphiquement.}},
  experimentales = {},
  numeriques = {},
}

\cpgeligne{MPSIA/A2/3/3}{
  parent = {MPSIA/A2/3},
  numero = {A2.3.3},
  contenu = {
    \cpgeitem{Développements limités.}},
  capacites = {
    \cpgeitem{Connaître et utiliser les développements limités à l’ordre 1 des fonctions $(1 + x)^\alpha$, $\mathrm{e}^x$ et $\ln(1 + x)$, et à l’ordre 2 des fonctions $\cos(x)$ et $\sin(x)$.}},
  experimentales = {},
  numeriques = {},
}

\cpgeligne{MPSIA/A2/3/4}{
  parent = {MPSIA/A2/3},
  numero = {A2.3.4},
  contenu = {
    \cpgeitem{Primitive et intégrale.}},
  capacites = {
    \cpgeitem{Interpréter l’intégrale comme une somme de contributions infinitésimales, en lien avec la méthode des rectangles en mathématiques.}},
  experimentales = {},
  numeriques = {},
}

\cpgeligne{MPSIA/A2/3/5}{
  parent = {MPSIA/A2/3},
  numero = {A2.3.5},
  contenu = {
    \cpgeitem{Valeur moyenne.}},
  capacites = {
    \cpgeitem{Exprimer la valeur moyenne sous forme d’une intégrale. Connaître la valeur moyenne sur une période des fonctions $\cos$, $\sin$, $\cos^2$ et $\sin^2$.}},
  experimentales = {},
  numeriques = {},
}

\cpgeligne{MPSIA/A2/3/6}{
  parent = {MPSIA/A2/3},
  numero = {A2.3.6},
  contenu = {
    \cpgeitem{Représentation graphique d’une fonction.}},
  capacites = {
    \cpgeitem{Déterminer un comportement asymptotique ; rechercher un extremum local.}
    \cpgeitem{Utiliser des échelles logarithmiques ; identifier une loi de puissance à une droite en échelle log-log.}},
  experimentales = {},
  numeriques = {},
}

\cpgeligne{MPSIA/A2/3/7}{
  parent = {MPSIA/A2/3},
  numero = {A2.3.7},
  contenu = {
    \cpgeitem{Développement en série de Fourier d’une fonction périodique.}},
  capacites = {
    \cpgeitem{Utiliser un développement en série de Fourier fourni par un formulaire.}},
  experimentales = {},
  numeriques = {},
}

\cpgerubrique{MPSIA/A2/4}{
  parent = {MPSIA/A2},
  numero = {A2.4},
  niveau = {2},
  titre = {Géométrie},
  chapeau = {},
}

\cpgeligne{MPSIA/A2/4/1}{
  parent = {MPSIA/A2/4},
  numero = {A2.4.1},
  contenu = {
    \cpgeitem{Vecteurs et système de coordonnées.}},
  capacites = {
    \cpgeitem{Exprimer les coordonnées d’un vecteur dans une base orthonormée.}
    \cpgeitem{Utiliser les systèmes de coordonnées cartésiennes, cylindriques et sphériques.}},
  experimentales = {},
  numeriques = {},
}

\cpgeligne{MPSIA/A2/4/2}{
  parent = {MPSIA/A2/4},
  numero = {A2.4.2},
  contenu = {
    \cpgeitem{Projection d'un vecteur et produit scalaire.}},
  capacites = {
    \cpgeitem{Interpréter géométriquement le produit scalaire et connaître son expression en fonction des coordonnées dans une base orthonormée.}
    \cpgeitem{Utiliser la bilinéarité et le caractère symétrique du produit scalaire.}},
  experimentales = {},
  numeriques = {},
}

\cpgeligne{MPSIA/A2/4/3}{
  parent = {MPSIA/A2/4},
  numero = {A2.4.3},
  contenu = {
    \cpgeitem{Produit vectoriel.}},
  capacites = {
    \cpgeitem{Interpréter géométriquement le produit vectoriel et connaître son expression en fonction des coordonnées dans une base orthonormée directe. Utiliser la bilinéarité et le caractère antisymétrique du produit vectoriel.}
    \cpgeitem{Faire le lien avec l’orientation des trièdres.}},
  experimentales = {},
  numeriques = {},
}

\cpgeligne{MPSIA/A2/4/4}{
  parent = {MPSIA/A2/4},
  numero = {A2.4.4},
  contenu = {
    \cpgeitem{Transformations géométriques.}},
  capacites = {
    \cpgeitem{Utiliser les symétries par rapport à un plan, les translations et les rotations de l’espace.}
    \cpgeitem{Utiliser leur effet sur l’orientation de l’espace.}},
  experimentales = {},
  numeriques = {},
}

\cpgeligne{MPSIA/A2/4/5}{
  parent = {MPSIA/A2/4},
  numero = {A2.4.5},
  contenu = {
    \cpgeitem{Courbes planes.}},
  capacites = {
    \cpgeitem{Reconnaître l’équation cartésienne d’une droite, d’un cercle.}
    \cpgeitem{Utiliser la représentation polaire d’une courbe plane ; utiliser un grapheur pour obtenir son tracé.}},
  experimentales = {},
  numeriques = {},
}

\cpgeligne{MPSIA/A2/4/6}{
  parent = {MPSIA/A2/4},
  numero = {A2.4.6},
  contenu = {
    \cpgeitem{Courbes planes paramétrées.}},
  capacites = {
    \cpgeitem{Identifier une ellipse à l’aide de sa représentation paramétrique $(x = a\cos(\omega t)$, $y = b\cos(\omega t-\varphi))$ et la tracer dans les cas particuliers $\varphi = 0$, $\varphi = \pi/2$ et}
    \cpgeitem{$\varphi = \pi$.}},
  experimentales = {},
  numeriques = {},
}

\cpgeligne{MPSIA/A2/4/7}{
  parent = {MPSIA/A2/4},
  numero = {A2.4.7},
  contenu = {
    \cpgeitem{Longueurs, aires et volumes classiques.}},
  capacites = {
    \cpgeitem{Citer les expressions du périmètre d’un cercle, de l’aire d’un disque, de l’aire d’une sphère, du volume d’une boule, du volume d’un cylindre.}},
  experimentales = {},
  numeriques = {},
}

\cpgeligne{MPSIA/A2/4/8}{
  parent = {MPSIA/A2/4},
  numero = {A2.4.8},
  contenu = {
    \cpgeitem{Barycentre d'un système de points.}},
  capacites = {
    \cpgeitem{Énoncer la définition du barycentre. Utiliser son associativité. Exploiter les symétries pour prévoir la position du barycentre d’un système homogène.}},
  experimentales = {},
  numeriques = {},
}

\cpgerubrique{MPSIA/A2/5}{
  parent = {MPSIA/A2},
  numero = {A2.5},
  niveau = {2},
  titre = {Trigonométrie},
  chapeau = {},
}

\cpgeligne{MPSIA/A2/5/1}{
  parent = {MPSIA/A2/5},
  numero = {A2.5.1},
  contenu = {
    \cpgeitem{Angle orienté.}},
  capacites = {
    \cpgeitem{Définir une convention d’orientation des angles d’un plan (euclidien) et lire des angles orientés. Relier l'orientation d'un axe de rotation à l'orientation positive des angles d’un plan perpendiculaire à cet axe.}},
  experimentales = {},
  numeriques = {},
}

\cpgeligne{MPSIA/A2/5/2}{
  parent = {MPSIA/A2/5},
  numero = {A2.5.2},
  contenu = {
    \cpgeitem{Fonctions cosinus, sinus et tangente.}},
  capacites = {
    \cpgeitem{Utiliser le cercle trigonométrique et l’interprétation géométrique des fonctions cosinus, sinus et tangente comme aide-mémoire : relation $\cos^2 x + \sin^2 x = 1$, relations entre fonctions trigonométriques et toutes relations du type $\cos(\pi \pm x)$ et $\cos(\pi/2 \pm x)$, parités, périodicité, valeurs des fonctions pour les angles usuels.}
    \cpgeitem{Citer les formules d’addition et de duplication des cosinus et sinus ; utiliser un formulaire dans les autres cas.}},
  experimentales = {},
  numeriques = {},
}

\cpgeligne{MPSIA/A2/5/3}{
  parent = {MPSIA/A2/5},
  numero = {A2.5.3},
  contenu = {
    \cpgeitem{Nombres complexes et représentation dans le plan. Somme et produit de nombres complexes.}},
  capacites = {
    \cpgeitem{Calculer et interpréter géométriquement la partie réelle, la partie imaginaire, le module et l’argument d’un nombre complexe.}},
  experimentales = {},
  numeriques = {},
}

\cpgerubrique{MPSIA/A2/6}{
  parent = {MPSIA/A2},
  numero = {A2.6},
  niveau = {2},
  titre = {Analyse vectorielle},
  chapeau = {},
}

\cpgeligne{MPSIA/A2/6/1}{
  parent = {MPSIA/A2/6},
  numero = {A2.6.1},
  contenu = {
    \cpgeitem{Gradient d’un champ scalaire.}},
  capacites = {
    \cpgeitem{Citer le lien entre le gradient et la différentielle. Citer l’expression de la différentielle en fonction des dérivées partielles.}
    \cpgeitem{Citer l’expression du gradient en coordonnées cartésiennes ; utiliser un formulaire fourni en coordonnées cylindriques ou sphériques.}
    \cpgeitem{Utiliser le fait que le gradient d’une fonction $f$ est perpendiculaire aux surfaces iso-$f$ et orienté dans le sens des valeurs de $f$ croissantes.}},
  experimentales = {},
  numeriques = {},
}

\cpgerubrique{MPSIA/A3}{
  parent = {MPSI},
  numero = {A3},
  niveau = {1},
  titre = {Outils numériques},
  chapeau = {
    % Le texte officiel écrit « la formation des étudiants inclue » : coquille corrigée
    % en « inclut » (voir la documentation, « Écarts au texte officiel »).
    \cpgeitem{La prise en compte de capacités de codage en langage Python dans la formation des étudiants inclut l’utilisation de fonctions extraites de diverses bibliothèques. Elle vise à une meilleure appréhension des principes mis en œuvre par les différents logiciels de traitement des données dont l’utilisation est par ailleurs toujours recommandée. Elle a aussi pour objectif de mobiliser ces capacités dans un contexte concret, celui de la physique et de la chimie. Cette formation par le codage permet également de développer des capacités utiles à la physique-chimie comme le raisonnement, la logique ou la décomposition d’un problème complexe en étapes plus simples.}
    \cpgeitem{Le tableau ci-dessous explicite ces outils ainsi que les capacités exigibles en fin de première année. Il sera complété dans le programme de physique-chimie de seconde année.}},
}

\cpgerubrique{MPSIA/A3/1}{
  parent = {MPSIA/A3},
  numero = {A3.1},
  niveau = {2},
  titre = {Outils graphiques},
  chapeau = {},
}

\cpgeligne{MPSIA/A3/1/1}{
  parent = {MPSIA/A3/1},
  numero = {A3.1.1},
  contenu = {
    \cpgeitem{Représentation graphique d’un nuage de points.}},
  capacites = {},
  experimentales = {},
  numeriques = {
    \cpgeitem{Utiliser les fonctions de base de la bibliothèque \texttt{matplotlib} pour représenter un nuage de points.}},
}

\cpgeligne{MPSIA/A3/1/2}{
  parent = {MPSIA/A3/1},
  numero = {A3.1.2},
  contenu = {
    \cpgeitem{Représentation graphique d’une fonction.}},
  capacites = {},
  experimentales = {},
  numeriques = {
    \cpgeitem{Utiliser les fonctions de base de la bibliothèque \texttt{matplotlib} pour tracer la courbe représentative d’une fonction.}},
}

\cpgeligne{MPSIA/A3/1/3}{
  parent = {MPSIA/A3/1},
  numero = {A3.1.3},
  contenu = {
    \cpgeitem{Courbes planes paramétrées.}},
  capacites = {},
  experimentales = {},
  numeriques = {
    \cpgeitem{Utiliser les fonctions de base de la bibliothèque \texttt{matplotlib} pour tracer une courbe plane paramétrée.}},
}

\cpgerubrique{MPSIA/A3/2}{
  parent = {MPSIA/A3},
  numero = {A3.2},
  niveau = {2},
  titre = {Équations algébriques},
  chapeau = {},
}

\cpgeligne{MPSIA/A3/2/1}{
  parent = {MPSIA/A3/2},
  numero = {A3.2.1},
  contenu = {
    \cpgeitem{Résolution d’une équation algébrique ou d’une équation transcendante : méthode dichotomique.}},
  capacites = {},
  experimentales = {},
  numeriques = {
    \cpgeitem{Déterminer, en s’appuyant sur une représentation graphique, un intervalle adapté à la recherche numérique d’une racine par une méthode dichotomique.}
    \cpgeitem{Mettre en œuvre une méthode dichotomique afin de résoudre une équation avec une précision donnée.}
    \cpgeitem{Utiliser la fonction \texttt{bisect} de la bibliothèque \texttt{scipy\.optimize} (sa spécification étant fournie).}},
}

\cpgerubrique{MPSIA/A3/3}{
  parent = {MPSIA/A3},
  numero = {A3.3},
  niveau = {2},
  titre = {Intégration - Dérivation},
  chapeau = {},
}

\cpgeligne{MPSIA/A3/3/1}{
  parent = {MPSIA/A3/3},
  numero = {A3.3.1},
  contenu = {
    \cpgeitem{Calcul approché d’une intégrale sur un segment par la méthode des rectangles.}},
  capacites = {},
  experimentales = {},
  numeriques = {
    \cpgeitem{Mettre en œuvre la méthode des rectangles pour calculer une valeur approchée d’une intégrale sur un segment.}},
}

\cpgeligne{MPSIA/A3/3/2}{
  parent = {MPSIA/A3/3},
  numero = {A3.3.2},
  contenu = {
    \cpgeitem{Calcul approché du nombre dérivé d’une fonction en un point.}},
  capacites = {},
  experimentales = {},
  numeriques = {
    \cpgeitem{Utiliser un schéma numérique pour déterminer une valeur approchée du nombre dérivé d’une fonction en un point.}},
}

\cpgerubrique{MPSIA/A3/4}{
  parent = {MPSIA/A3},
  numero = {A3.4},
  niveau = {2},
  titre = {Équations différentielles},
  chapeau = {},
}

\cpgeligne{MPSIA/A3/4/1}{
  parent = {MPSIA/A3/4},
  numero = {A3.4.1},
  contenu = {
    \cpgeitem{Équations différentielles d’ordre 1.}},
  capacites = {},
  experimentales = {},
  numeriques = {
    \cpgeitem{Mettre en œuvre la méthode d’Euler explicite afin de résoudre une équation différentielle d’ordre 1.}},
}

\cpgeligne{MPSIA/A3/4/2}{
  parent = {MPSIA/A3/4},
  numero = {A3.4.2},
  contenu = {
    \cpgeitem{Équations différentielles d’ordre supérieur ou égal à 2}},
  capacites = {},
  experimentales = {},
  numeriques = {
    \cpgeitem{Transformer une équation différentielle d’ordre n en un système différentiel de n équations d’ordre 1.}
    \cpgeitem{Utiliser la fonction \texttt{odeint} de la bibliothèque \texttt{scipy\.integrate} (sa spécification étant fournie).}},
}

\cpgerubrique{MPSIA/A3/5}{
  parent = {MPSIA/A3},
  numero = {A3.5},
  niveau = {2},
  titre = {Probabilités - statistiques},
  chapeau = {},
}

\cpgeligne{MPSIA/A3/5/1}{
  parent = {MPSIA/A3/5},
  numero = {A3.5.1},
  contenu = {
    \cpgeitem{Variable aléatoire.}},
  capacites = {},
  experimentales = {},
  numeriques = {
    \cpgeitem{Utiliser les fonctions de base des bibliothèques \texttt{random} et/ou \texttt{numpy} (leurs spécifications étant fournies) pour réaliser des tirages d’une variable aléatoire.}
    \cpgeitem{Utiliser la fonction \texttt{hist} de la bibliothèque \texttt{matplotlib.pyplot} (sa spécification étant fournie) pour représenter les résultats d’un ensemble de tirages d’une variable aléatoire.}
    \cpgeitem{Déterminer la moyenne et l’écart-type d’un ensemble de tirages d’une variable aléatoire.}},
}

\cpgeligne{MPSIA/A3/5/2}{
  parent = {MPSIA/A3/5},
  numero = {A3.5.2},
  contenu = {
    \cpgeitem{Régression linéaire.}},
  capacites = {},
  experimentales = {},
  numeriques = {
    \cpgeitem{Utiliser la fonction \texttt{polyfit} de la bibliothèque \texttt{numpy} (sa spécification étant fournie) pour exploiter des données.}
    \cpgeitem{Utiliser la fonction \texttt{random\.normal} de la bibliothèque \texttt{numpy} (sa spécification étant fournie) pour simuler un processus aléatoire.}},
}

